\documentclass{article}

% if you need to pass options to natbib, use, e.g.:
% \PassOptionsToPackage{numbers, compress}{natbib}
% before loading nips_2018

% ready for submission
%\usepackage{nips_2018} % For anonymous submission %%PA.5.17: we need this one for submission (otherwise we get automatically rejected...)
%\usepackage[preprint]{nips_2018} % For arxiv
%\usepackage[final]{nips_2018} % For final submission

% to compile a preprint version, e.g., for submission to arXiv, add
% add the [preprint] option:
% \usepackage[preprint]{nips_2018}

% to compile a camera-ready version, add the [final] option, e.g.:
\usepackage[preprint]{nips_2018}

% to avoid loading the natbib package, add option nonatbib:
% \usepackage[nonatbib]{nips_2018}

\usepackage{rotating} % for rotating the tables to fit in

\usepackage[utf8]{inputenc} % allow utf-8 input
\usepackage[T1]{fontenc}    % use 8-bit T1 fonts
\usepackage{hyperref}       % hyperlinks
\usepackage{url}            % simple URL typesetting
\usepackage{booktabs}       % professional-quality tables
\usepackage{amsfonts}       % blackboard math symbols
\usepackage{nicefrac}       % compact symbols for 1/2, etc.
\usepackage{microtype}      % microtypography
%\usepackage{natbib}
\usepackage{amsmath}
\usepackage{algorithm}
\usepackage{algpseudocode}
%\usepackage{algorithm2e}
%\usepackage{algorithmic}
\usepackage{float}% http://ctan.org/pkg/float
\usepackage{graphicx}
\usepackage{subfig}
\usepackage{placeins}

\allowdisplaybreaks[4]

% ----------------------------------------------------------------------------------------
% algorithmic environment
%\usepackage[noend]{algpseudocode}
\usepackage{placeins}
%\usepackage{showlabels}
\usepackage{algorithm}
%\usepackage{algpseudocode}

\usepackage{etoolbox}
\usepackage{xr}
\externaldocument{../X_learner_PNAS_main/X-Learner_PNAS-main}

\makeatletter  % needed for pseudocode
\def\BState{\State\hskip-\ALG@thistlm}
\makeatother

\makeatletter
\AfterEndEnvironment{algorithm}{\let\@algcomment\relax}
\AtEndEnvironment{algorithm}{\kern2pt\hrule\relax\vskip3pt\@algcomment}
\let\@algcomment\relax
\newcommand\algcomment[1]{\def\@algcomment{\footnotesize#1}}

\renewcommand\fs@ruled{\def\@fs@cfont{\bfseries}\let\@fs@capt\floatc@ruled
	\def\@fs@pre{\hrule height.8pt depth0pt \kern2pt}%
	\def\@fs@post{}%
	\def\@fs@mid{\kern2pt\hrule\kern2pt}%
	\let\@fs@iftopcapt\iftrue}
\makeatother
% ----------------------------------------------------------------------------------------

\newtheorem{assumption}{Assumption}

\title{Transfer Learning for Estimating Causal Effects using Neural Networks}
% CAUSAL NEURAL NETWORKS means something pretty different to most of the AI community. 
% So we should find a way to word things that is not too sensationalist. 


% The \author macro works with any number of authors. There are two
% commands used to separate the names and addresses of multiple
% authors: \And and \AND.
%
% Using \And between authors leaves it to LaTeX to determine where to
% break the lines. Using \AND forces a line break at that point. So,
% if LaTeX puts 3 of 4 authors names on the first line, and the last
% on the second line, try using \AND instead of \And before the third
% author name.


\newcommand*\samethanks[1][\value{footnote}]{\footnotemark[#1]}
\author{
	Sören R. Künzel\thanks{These authors contributed equally to this work.} \\
%	Statistics\\
	UC Berkeley\\
	\texttt{srk@berkeley.edu} \\
	\And
	Bradly C. Stadie\samethanks \\
%	Statistics\\
	UC Berkeley\\
	\texttt{bstadie@berkeley.edu} \\
	\And	
	Nikita Vemuri \\
%	EECS\\
	UC Berkeley \\
%	\texttt{nikitavemuri@berkeley.edu } \\
	\AND
	Varsha Ramakrishnan \\
%	EECS\\
	UC Berkeley \\
%	\texttt{vio@berkeley.edu} \\
	\And
	Jasjeet S. Sekhon \\ % BCS: JAS AND I DISCUSSED AND AGREED ON THIS 							        ORDER
%	Poli.\ Sci. and Statistics\\
	UC Berkeley \\
	\texttt{sekhon@berkeley.edu} \\
	\And
    Pieter Abbeel \\
%	EECS\\
	UC Berkeley \\
	\texttt{pabbeel@cs.berkeley.edu} \\
}

\begin{document}
	\maketitle
	\begin{abstract}
		We develop new algorithms for estimating heterogeneous treatment effects, combining recent developments in transfer learning for neural networks with insights from the causal inference literature. By taking advantage of transfer learning, we are able to efficiently use different data sources that are related to the same underlying causal mechanisms. We compare our algorithms with those in the extant literature using extensive simulation studies based on large-scale voter persuasion experiments and the MNIST database. Our methods can perform an order of magnitude better than existing benchmarks while using a fraction of the data.
	\end{abstract}
	
	\section{Introduction} \label{section:Introduction}

The rise of massive datasets that provide fine-grained information about human beings and their behavior provides unprecedented opportunities for evaluating the effectiveness of treatments. Researchers want to exploit these large and heterogeneous datasets, and they often seek to estimate how well a given treatment works for individuals conditioning on their observed covariates. This problem is important in medicine (where it is sometimes called personalized medicine) \citep{Henderson2016,powers2018some}, digital experiments \citep{taddy2016nonparametric}, economics \citep{Athey2016}, political science \citep{green2012modeling}, statistics
\citep{tian2014simple}, and many other fields. A large number of articles are being written on this topic, but many outstanding questions remain. We present the first paper that applies transfer learning to this problem.

%The rise of massive datasets that provide fine-grained information about human beings and their behavior provides unprecedented opportunities for evaluating the effectiveness of treatments in a wide-variety of fields ranging from clinical medicine to political science to online digital experiments. Given the large and heterogeneous datasets, researchers and practitioners are increasingly unsatisfied with estimating how well treatments work on \textit{average}. Instead, they often seek to estimate how well a given treatment works for individuals conditioning on their observed covariates. This allows them to personalize treatment recommendations.

% Classical only one experiment, but a lot of information left on the table. It is different from prediction, and one has to leverage information already. We need to use this information
In the simplest case, treatment effects are estimated by splitting a training set into a treatment and a control group. The treatment group receives the treatment, while the control group does not. The outcomes in those groups are then used to construct an estimator for the Conditional Average Treatment Effect (CATE), which is defined as the expected outcome under treatment minus the expected outcome under control given a particular feature vector \citep{athey2015machine}. This is a challenging task because, for every unit, we either observe its outcome under treatment or control, but never both. Assumptions, such as the random assignment of treatment and additional regularity conditions, are needed to make progress. Even with these assumptions, the resulting estimates are often noisy and unstable because the CATE is a vector parameter. Recent research has shown that it is important to use estimators which consider both treatment groups simultaneously %and leverage information between the treated and control groups 
(\cite{kunzel2017meta, wager2015estimation, nie2017learning, hill2011bayesian}). Unfortunately, these recent advances are often still insufficient to train robust CATE estimators because of the large sample sizes required when the number of covariates is not small.

%Introduce the idea of transfer learning. 
However, researchers usually fail to use ancillary datasets that are available to them in applications. This is surprising, given the need for additional data to estimate CATE reliably. These ancillary datasets are related to the causal mechanism under investigation, but they are also partially distinct so they cannot be pooled naively, which explains why researches often do not use them. Examples of such ancillary datasets include observations from: experiments in different locations on different populations, different treatment arms, different outcomes, and non-experimental observational studies. The key idea underlying our contributions is that one can substantially improve CATE estimators by transferring information from other data sources.

\textbf{Our contributions are as follows:}

\begin{enumerate}

\item \textbf{We introduce the new problem of transfer learning for estimating heterogeneous treatment effects.}

\item \textbf{We develop the Y-learner for CATE estimation.} We consider the problem of CATE estimation with deep neural networks. We propose the Y-Learner, a CATE estimator designed from the ground up to take advantage of deep neural networks' ability to easily share information across layers. The Y-Learner often achieves state-of-the-art performance on CATE estimation. The Y-learner does not use transfer learning.

%\item \textbf{We introduce the new problem of transfer learning for CATE estimation and propose several solutions.} \textbf{We show that it works in a difficult data problem.}

%\item \textbf{We adapt the Reptile transfer learning method to CATE estimation}. It is not immediately obvious how one should use Reptile for CATE estimation, as CATE estimation usually requires simultaneously estimating two quantities of interest and Reptile typically adapts on a per-task basis. We provide multiple sets of instructions for overcoming these difficulties. In adapting Reptile to our problem, we discovered a slight modification to the original algorithm. To distinguish this modification, we refer to it in this paper as \textbf{SF Reptile}, Slow-Fast Reptile. 

\item \textbf{MLRW Transfer for CATE Estimation} adapts the idea of meta-learning regression weights (MLRW) to CATE estimation. Using these learned weights, regression problems can be optimized much more quickly than with random initializations. Though a variety of MLRW algorithms exist, it is not immediately obvious how one should use these methods for CATE estimation. The principle difficulty is that CATE estimation requires the simultaneous estimation of outcomes under both treatment and control, when we only observe one of the outcomes for any individual unit. However, most MLRW transfer methods optimize on a per-task basis to estimate a single quantity. We show that one can overcome this problem with clever use of the Reptile algorithm \citep{reptile}.%%PA.5.17: this citation is off, should be Nichol, Achiam, Schulman, 2018?
While adapting Reptile to work with our problem, we discovered a slight modification to the original algorithm. To distinguish this modification, we refer to it in this paper as \textbf{SF Reptile}, Slow-Fast Reptile. 

\item \textbf{We provide several additional methods for transfer learning for CATE estimation:} warm start, frozen-features, multi-head, and joint training.

\item \textbf{We apply our methods to difficult data problems and show that they perform better than existing benchmarks.} We reanalyze a set of large field experiments that evaluate the effect of a mailer on voter turnout in the 2014 U.S. midterm elections \citep{gerber2017generalizability}. This includes 17 experiments with 1.96 million individuals in total. We also simulate several randomized controlled trials using image data of handwritten digits found in the MNIST database \citep{lecun1998mnist}. We show that our methods, \textbf{MLRW} in particular, obtain better than state-of-the-art performance in estimating CATE, and that they require far fewer observations than extant methods.

\item \textbf{We provide open source code for our algorithms.}\footnote{The software will be released once anonymity is no longer needed. We can also provide an anynomized copy to reviewers upon request.} 
\end{enumerate}

%We transfer learn to CATE estimation. Though a variety of methods for MLIW transfer exist, it is not immediately obvious how one should use these methods for CATE estimation. The principle difficulty is CATE estimation requires the simultaneous estimation of both outcome under treatment and outcome under control. However, most MLIW transfer methods adapt on a per-task basis to estimate a single quantity. We show that one can overcome this problem by adapting the MLIW algorithm Reptile. In adapting Reptile to our problem, we discovered a slight modification to the original algorithm. To distinguish this modification, we refer to it in this paper as \textbf{SF Reptile}, Slow-Fast Reptile.

%Statement of contributions. 
%In this paper, we consider the problem of transfer learning for CATE estimation. We will be particularly interested in transfer across different CATE estimation experiments that share similar units. This is a well studied and perennial problem in the literature, spanning several fields such as computer vision \cite{tuning}, semi-supervised learning \cite{zemmel, prognets, bunchofstuff}, and reinforcement learning \cite{MAML, Reptile}. Recent advances in transfer learning have typically exploited the structure of deep neural networks. Therefore, we will first need to consider the problem of CATE estimation with deep neural networks. This will lead to the development of the Y-Learner, a CATE estimator designed from the ground up to take advantage of deep neural network's ability to easily share information across layers. The Y-Learner often achieves state-of-the art performance on CATE estimation. Subsequently, we will investigate a variety of techniques for transfer learning across CATE estimators. Several of the techniques such as fine tuning, using previously trained features, and using networks with shared base layers and multiple heads are well-studied in neural network literature but new to this problem. Other techniques, such as SF Reptile, are novel contributions to the transfer learning literature.  

% We use the self-treatment to show how well it works.
%Our lead example aims to estimate the treatment effect of a mailer which is sent to potential voters to increase the election turnout for the primary election in Michigan 2006 \cite{GerberGreenLarimer}.
%The sample size is here with 229,461 units relatively large, but CATE estimation still turns out to be difficult.  
%In this case, we want to use two different data sets to improve the CATE estimation. First of all, in the same election, three other mailers have been tested to increase the voter turnout for the same election. Those mailers are different, but we still believe that there are communalities in the mechanism how those mailers work. 
%The mailer of interest is very popular, and we have found seven studies where this mailer was used. In all of these studies, the election was different from the one we are interested in, but the mailer was almost identical to the one we are interested in. 
%We will show how both additional data sources can significantly improve the final CATE estimator. 

%While our lead example seems to be very special for voter persuasion, similar situations exist on a much larger scale when targeting online advertisements with thousands of experiments, medical studies where the sample sizes are generally very small, and one has to use external data sources to construct good CATE estimators. To assuage any concerns that our methods are over-fit to this example, results are also presented on an example using MNIST digits extended from \cite{Yu2013stability}. Code for this project is available to reviewers upon request, and will be fully released along with the conference version of this paper.  %BCS 05/12: potentially MNIST Stuff. Do want to ease fears that we are not over-fitting to one problem. 

%BCS: Introduction may be long enough already. 
%The contributions of this paper are as follows 
%\begin{enumerate}
%	\item jjj
%\end{enumerate}

% Contributions of our paper (Neural Network Learners, Y Learner, Transfer Learning for Causal Inference, SF Reptile) and the setup

%BCS: Paragraphs like the below tend not to work well for NIPS submissions. The reviewers have very small attention spans and prefer that the contributions are explicitly listed. What I'm saying is the below paragraph will not even be read by 90 percent of reviewers. They need the contributions more explicitly stated in a list format. 

%In this paper, we start by introducing the problem formally and applying ideas of CATE estimation which already exist in the literature to neural networks. Here we introduce a new CATE estimator, the Y-Learner which performs extremely well for our data set. 
%Then we discuss different transfer learning strategies which have widely been used in the deep learning community, and we show how those ideas can use outside data sources to significantly improve the CATE estimators. 
%We then compare our estimators in an exhaustive simulation study before we apply them to real data. 

	
	\newcommand{\Pcal}{\mathcal{P}}
\newcommand{\R}{\mathbb{R}}
\newcommand{\Yobs}{Y^{obs}}
\newcommand{\E}{\mathbb{E}}
\newcommand{\define}{:=}
\renewcommand{\P}{\mathbb{P}}
\newcommand{\emin}{e_{\mbox{\small min}}}
\newcommand{\emax}{e_{\mbox{\small max}}}

\section{CATE ESTIMATION}

%\subsection{Background and Notation}
We begin by formally introducing the CATE estimation problem. 
Following the potential outcomes framework \citep{rubin1974estimating}, assume there exists a single experiment wherein we observe $N$ i.i.d. distributed units from some super population, $(Y_i(0), Y_i(1), X_i, W_i) \sim \Pcal$. $Y_i(0) \in \R$ denotes the potential outcome of unit $i$ if it is in the control group, $Y_i(1)\in \R$ is the potential outcome of $i$ if it is in the treatment group, $X_i \in \R^d$ is a $d$-dimensional feature vector, and $W_i \in \{0,1\}$ is the treatment assignment.
For each unit in the treatment group ($W_i = 1$), we only observe the outcome under treatment, $Y_i(1)$. For each unit under control ($W_i = 0$), we only observe the outcome under control. Crucially, there cannot exist overlap between the set of units for which $W_i = 1$ and the set for which $W_i = 0$. It is impossible to observe both potential outcomes for any unit. This is commonly referred to as the fundamental problem of causal inference. 


However, not all hope is lost. We can still estimate the Conditional Average Treatment Effect (CATE) of the treatment. Let $x$ be an individual feature vector. Then the CATE of $x$, denoted $\tau(x)$, is defined by %Let $\Yobs_i = Y_i(W_i)$ be the observed outcome. 

$$\tau(x) = \E[Y(1) - Y(0) | X = x].$$

Estimating $\tau$ is impossible without making further assumptions on the distribution of $(Y_i(0), Y_i(1), X_i, W_i)$. In particular, we need to place two assumptions on our data. 


%To avoid this problem, we assume that the treatment assignment is independent of the potential outcomes given the observed feature vector. Nevertheless, we can make CATE estimation tractable with two key assumptions.  .

%The main task is to use the observed data $D = (\Yobs_i, X_i, W_i)_{i=1}^N$ to estimate the Conditional Average Treatment Effect (CATE) for a new unit with feature vector $x$, 
%$$\tau(x) = \E[Y(1) - Y(0) | X = x].$$
%This is a difficult problem, and without making further assumptions on the distribution of $(Y_i(0), Y_i(1), X_i, W_i)$ it is not possible to estimate it. 
%For example, if there was some unobserved confounder, a random variable which influences both the probability of treatment and the potential outcomes, then the CATE is unidentifiable. To avoid this problem, we assume that the treatment assignment is independent of the potential outcomes given the observed feature vector.
\begin{assumption}[Strong Ignorability, \cite{rosenbaum1983central}] \label{assumption:StrongIgnorability}
$$
    (Y_i(1), Y_i(0)) \perp W |X.     
$$
\end{assumption}
\begin{assumption}[Overlap] \label{assumption:Overlap}
Define the propensity score of $x$ as, 
$$e(x) \define \P(W = 1 | X = x).$$ 
Then there exists constant $0< \emin, ~\emax < 1$ such that for all $x \in \mbox{Support}(X)$, 
$$0 < \emin < e(x) < \emax < 1.$$
In words, $e(x)$ is bounded away from 0 and 1.
\end{assumption}
Assumption \ref{assumption:StrongIgnorability} ensures that there is no unobserved confounder, a random variable which influences both the probability of treatment and the potential outcomes, which would make the CATE unidentifiable. The assumption is particularly strong and difficult to check in applications. Meanwhile, Assumption \ref{assumption:Overlap} rectifies the situation wherein a certain part of the population is always treated or always in the control group. If, for example, all women were in the control group, one cannot identify the treatment effect for women. Though both assumptions are strong, they are nevertheless satisfied by design in randomized controlled trials. While the estimators we discuss would be sensible in observational studies when the assumptions are satisfied, we warn practitioners to be cautious in such studies, especially when the number of covariates is large \citep{overlapAlex}.

Given these two assumptions, there exist many valid CATE estimators. The crux of these methods is to estimate two quantities: 
the control response function, 
$$
\mu_0(x) = \E[Y(0)|X=x],
$$
and the treatment response function,
$$
\mu_1(x) = \E[Y(1)|X=x].
$$
If we denote our learned estimates as $\hat\mu_0(x)$ and $\hat\mu_1(x)$, then we can form the CATE estimate as the difference between the two
$$
\hat \tau(x) = \hat \mu_1(x) - \hat \mu_0(x).
$$
The astute reader may be wondering why we don't simply estimate $\mu_0$ and $\mu_1$ with our favorite function approximation algorithm at this point and then all go home. After all, we have access to the ground truths $\mu_0$ and $\mu_1$ and the corresponding inputs $x$. In fact, it is commonplace to do exactly that. When people directly estimate $\mu_0$ and $\mu_1$ with their favorite model, we call the procedure a T-learner \citep{kunzel2017meta}. Common choices of models include linear models and random forests, though neural networks have recently been considered \citep{nie2017learning}. 

While it may seem like we've triumphed, the T-learner does have some drawbacks \citep{athey2015machine}. It is usually an inefficient estimator. For example, it will often perform poorly when one can borrow information across the treatment conditions. To overcome these deficiencies, a variety of alternative learners have been suggested. Closely related to the T-learner is the idea of estimating the outcome using all of the features and the treatment indicator, without giving the treatment indicator a
special role \citep{hill2011bayesian}. The predicted CATE for an individual unit is then the difference between the predicted values when the treatment assignment indicator is changed from control to treatment, with all other features held fixed. This is called the S-learner, because it uses a single prediction model.

In this paper, we suggest another new learner called the $\mathbf{Y}$\textbf{-learner} (See Figure \ref{fig:ylearner}). This learner has been engineered from the ground up to take advantage of some of the unique capabilities of neural networks. See the appendix for a full description of the Y-learner, and additional learners found in the literature. Below, we will use these learners as base algorithms for transfer learning. That is to say, we will use the knowledge gained by training one learner on one experiment to help a new learner with a new underlying experiment train faster with less data. 

\FloatBarrier
\begin{figure}
	\centering
	\includegraphics[width=0.5\linewidth]{figures/y_learner}
    \caption{Y-learner with Neural Networks. One of many advanced methods for CATE estimation. See the appendix Section \ref{section:XvsY} and \ref{app:transfer.algos} for a more detailed overview.}
    \label{fig:ylearner}
\end{figure}

%We follow \cite{kunzel2017meta} to categorize CATE estimators into CATE meta-learners. A CATE meta-learner decomposes the problem of estimating the CATE into several sub-regression problems, which can be solved using any regression estimator.\footnote{Note that the term "CATE meta learner" is different from the term "meta learners" used for neural networks, which usually describe transfer learning strategies for neural networks.


%The simplest example of a CATE meta-learner is the T-learner, which decomposes the CATE estimation procedure into the two separate sub-regression problems: estimating the control response function, $\mu_0(x) = \E[Y(0)|X=x]$, and the treatment response function, $\mu_1(x) = \E[Y(1)|X=x]$. The CATE is then estimated by subtracting the two estimates $\hat \tau(x) = \hat \mu_1(x) - \hat \mu_0(x)$.
%Note that $\mu_0$ and $\mu_1$ can be estimated using any regression method. For example, the following base learners have been used in the literature:  penalized linear models \citep{foster2013subgroup}, Random Forests (RF) \citep{athey2015machine}, and Bayesian Additive Regression Trees (BART) \citep{hill2011bayesian,green2012modeling}. 
%We implemented the T-learner with RF and Neural Networks (NN) as base learners, and we refer to these learners as T-RF and T-NN. 

%Similar to the T-learner, \cite{kunzel2017meta} identify four more meta-learners: the X, S, F, and the U learners. We evaluate all of these CATE meta-learners with neural networks and random forests as base learners. We use the notation X-NN to refer to the X-learner (which is formally defined below) that uses NN as the base learner. Therefore, we create a large pool of benchmark methods, which achieve state-of-the-art performance. 

%and to demonstrate that neural networks often outperform estimators which are based on trees such as RF and BART. We do not believe that neural networks generally perform better than RF, but in our evaluation section, we find that they work much better considering the characteristics of our datasets.

%R-NN, is out for now. We need to put this in later.
%Not all CATE estimators can be classified as meta-learners. 
%To our knowledge, the only previous CATE estimator using NN is \cite{nie2017learning}, and we refer to it as R-NN. It is similar to meta-learners in that it decomposes the problem into sub-problems which can be treated independently. Specifically, in a first step it estimates the propensity score and the mean outcome function, $\mu(x) = \E[\Yobs | X = x]$, and in a second step, it estimates the CATE by minimizing a new loss function which depends on the estimate of the propensity score and the estimate of the mean outcome function. We include it as a benchmark method.

%Neural Networks have proven to be extremely successful on many data sets, and as we will see in Section \ref{section:Simulation}, they perform well in our examples. Using them as base learners enables us to create very powerful CATE estimators which perform favorably compared to other baseline estimators, which can even handle images.
%Most importantly, neural networks enable us to borrow strength across different datasets using transfer learning.



%\subsection{Transfer Learning Background}
%The key idea of transfer learning is that new experiments should transfer knowledge from past experiences and external datasets. Perhaps the most straightforward example of successful transfer comes from computer vision. Here, large neural networks are trained to label the content of an image. Subsequently, the weights from this network are used to initialize a new neural network for a new task. 


%Suppose that a deep neural network $\pi_\theta(I_j) = \ell_j$ is trained to take an input image $I_j$ and output a label $\ell_j$ corresponding to the object depicted in the image. To classify thousands of images correctly, $\pi_\theta$ must learn to encode sufficiently general and useful features with $\pi_\theta$. In fact, we often find that these features are sufficiently general and useful that they are helpful.    

\section{Transfer Learning}


\subsection{Background}
The key idea in transfer learning is that new experiments should transfer insights from previous experiments rather than starting learning anew. The most straightforward example of transfer comes from computer vision \citep{finetune0, finetune2, finetune1, donahu}. %%PA.5.17: Donahue et al is a great example of transfer in NNs, but there is a lot of transfer work preceding it; It'd be good to first discuss in a couple sentences how there is a long history of research in transfer, cite a bunch of the oldest papers (Donahue et al could provide the references for this), and then say that our work builds most directly on / is most directly inspired by Donahue et al ...
Here, it is standard practice to train a neural network $\pi_\theta$ for one task and then use the trained network weights $\theta$ as initialization for a new task. The hope is that some basic low-level features of a vision system should be quite general and reusable. Starting optimization from networks that have already learned these general features should be faster than starting from scratch. 

Despite its promise, fine-tuning often fails to produce initializations that are uniformly good for solving new tasks \citep{maml}. One potent fix to this problem is a class of algorithms that seek to optimize meta-learning initialization weights \citep{maml, reptile}. In these algorithms, one meta-optimizes over many experiments to obtain neural network weights that can quickly find solutions to new experiments. We will use the Reptile algorithm to learn initialization weights for CATE estimation.   


%A recently proposed solution to this problem is to meta learn initialization weights that can optimize the loss function much more quickly than random initialization \cite{maml, reptile}. 


%By good initializers, we mean that starting from these weights, one can train neural networks $\pi_{\theta_0}$ and $\pi_{\theta_1}$ to estimate $\mu^i_0$ and $\mu^i_1$ for any arbitrary experiment much faster and with less data than starting from random initializations 


%Luckily, a large body of transfer learning literature exists to rectify these issues \cite{tonsofstuff}. In this paper, we adapt two such strategies: joint-training and reptile \cite{reptile}. These methods were chosen because they are relatively simple, flexible, and delivered good results. 



%It may be the case that $\theta^i_0 \cap \theta^i_1 = \emptyset$. However, they may also share layers as in the Y-Learner.

%For example, a vision system that is good at labeling animals in images should also learn features that are useful for segmenting parts of a photograph

%Sometimes, a slight modification of this idea is used wherein only some of the network weights will be used as `base layers' that feed their features into new untrained task-specific layers. In both cases, t


%More concretely, let $\theta = \left[\theta^1, \theta^2, \dots, \theta^N \right]$, with $\theta^i$ being the weights of the $i$th neural network layer. Then, a new task considers a mixed set of weights $\left[ \theta^1, \theta_2, \phi^1, \phi^2, \dots, \phi^M \right]$, where $\theta^i$ are transfered from a previous task and $\phi^i$ are freshly initialized for the current task (See figure??)


%manage to succeed when more naive methods fail because they explicitly optimize for finding a good initialization $\theta$ that can solve new tasks more efficiently than random initializations.


%Indeed, we found that in practice fine-tuning offered a minimal performance boost over vanilla estimators whereas joint-training and SF Reptile offered more substaintial gains. 

%The term joint-training is sometimes underspecified in the literature. In this paper, joint-training will be defined as follows. Suppose there exist $k$ tasks you would like to solve, each sharing a common input space. In joint-training, you train a neural network with base features $[\theta_1, \dots, \theta_i]$ and $k$ branches $[\theta_1, \dots, \theta_i, \phi_1^1, \dots, \phi_M^1],\dots, [\theta_1, \dots, \theta_i, \phi_1^k, \dots, \phi_M^k]$. Each branch is simultaneously trained and the gradients for the shared base layers are averaged. This interpretation is shown in figure ??. Meanwhile, in Reptile only one set of weights $[\theta_1, \dots , \theta_N]$ exists. There exist no task-specific weights. Instead, we train one set of weights $\theta$ to act as a good initialization for solving new tasks. By a good initialization, we mean that task $k+1$ can be solved with an order of magnitude less data when initialized from $\theta$ instead of a standard random initialization. Note that this is similar in spirit to the fine-tuning discussed above. However, in Reptile, we also optimize over the objective of $\theta$ being a good initialization by using a type of shrinkage optimizer. See our discussion of SF Reptile below for further details. Finally, a variant of MAML called first-order MAML is quite similar to reptile. Due to these similarities, we omit any further MAML analysis from this work.  

\subsection{Transfer Learning CATE Estimators}

%We are interested in transferring features from one average treatment effect estimate to another. Of the six methods discussed below, fine-tuning, frozen-features, multi-head, and joint training are relatively standard ideas in the field of neural networks but hitherto unexplored in the context of CATE estimation. Likewise, MLIW methods have been studied in the context of meta-reinforcement learning and supervised learning but no such algorithms exist for CATE estimation.  

%Assuming there is mutual information between the experiments, then positive transfer should enable faster learning with less data

%. Learning meta-regression weights that make good initializers has been studied in the context of 

%Using one set of meta-regression weights for CATE estimation is an entirely novel concept, as all previous CATE estimators, we are aware of use at least some distinct weights for estimating $\mu_0$ and $\mu_1$. 

%Though the discussion below typically assumes transfer between only two experiments; this is for ease of exposition. The ideas presented easily generalizes to transfer between an arbitrary number of experiments. In the next section, the problems we consider all involve transfer between $10-30$ experiments. Furthermore, in the case of the Y and X learners, $\pi_{\theta_i}$ requires computation from intermediate networks which must also be trained. For simplicity and so the discussion remains largely agnostic towards one's choice of CATE estimator, we omit those intermediate steps from the discussion below. We refer the reader to Appendix A for step-by-step algorithms for transfer learning with the T and Y learners. An even more granular level of detail is available in the released code. 

In this section, we consider a scenario wherein one has access to many related causal inference experiments. Across all experiments, the input space $X$ is the same. Let $i$ index an experiment. Each experiment has its own distinct outcome when treatment is received, $\mu^i_1(x)$, and when no treatment is received, $\mu^i_0(x)$. Together, these quantities define the CATE $\tau^i = \mu^i_1 - \mu^i_0$, which we want to estimate. We are usually interested in estimating the CATE by using $X$ to predict $\mu_0^i$ and $\mu_1^i$. However, in transfer learning, the hope is that we can transfer knowledge between experiments such that being able to predict $\mu_0^i, \mu_1^i,$ and $\tau^i$ from experiment $i$ accurately will help us predict $\mu_0^j, \mu_1^j$, and $\tau^j$ from experiment $j$. 

Below, let $\pi_{\theta}$ be a generic expression for a neural network parameterized by $\theta$. Sometimes, parameters will have a subscript indicating if their neural network predicts treatment or control (0 for control and 1 for treatment). Parameters may also have a superscript indicating the experiment number whose outcome is being predicted. For example, $\pi_{\theta_0^2}(x)$ predicts $\mu^2_0(x)$, the outcome under control for Experiment $2$. We will sometimes drop the superscript $i$ when the meaning is clear. All of the transfer algorithms described here are presented in detail in Appendix \ref{app:transfer.algos}.

\textbf{Warm start (also known as fine-tuning):} Experiment $0$ predicts $\pi_{\theta^0_0}(x) = \hat{\mu}^0_0(x)$ and $\pi_{\theta^0_1}(x) = \hat{\mu}^0_1(x)$ to form the CATE estimator $\hat{\tau} = \hat{\mu}^0_1(x) - \hat{\mu}^1_0(x)$. Suppose $\theta^0_0$, $\theta^0_1$ are fully trained and produce a good CATE estimate. For experiment 1, the input space $X$ is identical to the input space for experiment $0$, but the outcomes $\mu^1_0(x)$ and $\mu^1_1(x)$ are different. %%PA.5.17: I am confused, isn't for experiment 1 the output space described by D instead of mu?  I might be wrong, but either way, I think a lot of people who haven't studied CATE estimation before would be lost here, requires clarification...  maybe somewhere explicitly define what it means to be experiment 0 and what it means to be experiment 1?  then later explain each version; also, especially b/c it's unfamiliar terminology, it becomes hard when above things are called two step procedure, but here it's seemingly a two experiment procedure; consistency in terminology helps a lot when reading new material [or maybe these aren't referring to same two things, not sure, lost on that...]
However, we suspect the underlying data representations learned by $\pi_{\theta^0_0}$ and $\pi_{\theta^0_1}$ are still useful. Hence, rather than randomly initialize $\theta_0^1$ and $\theta_1^1$ for experiment 1, we set $\theta_0^1 = \theta_0^0$ and $\theta_1^1 = \theta_1^0$. We then train $\pi_{\theta_0^1}(x) = \hat{\mu}^1_0(x)$ and $\pi_{\theta_1^1}(x) = \hat{\mu}^1_1(x)$. See Figure \ref{fig:warm-start} and Algorithm \ref{alg:warm-t} in the appendix.

%\begin{figure}
%    \label{fig:warm-start}
%    \begin{center}
%            \includegraphics[scale=0.3]{figures/Warm_Start}
%        \caption{Warm Start for the T-learner}
%    \end{center}
%\end{figure}

\begin{figure}[t]
\begin{center}
\includegraphics[width = 1\textwidth]{figures/all_methods1}
\end{center}
%\begin{center}
%\subfloat{\includegraphics[width=0.30\textwidth]{figures/Warm_Start}}\hfill
%\subfloat{\includegraphics[width=0.30\textwidth]{figures/Frozen_Features}} \hfill
%\subfloat{\includegraphics[width=0.30\textwidth]{figures/Warm_Start}}
%\end{center}
\caption{Warm start, frozen-features, and multi-head methods for CATE transfer learning. For these figures, we use the T-learner as the base learner for simplicity. All three methods attempt to reuse neural network features from previous experiments. See the appendix for an illustration of joint-training.}
\label{fig:warm-start}
\end{figure} 

%Instead of initializing $\phi^1, \phi^2$ randomly, we set $\phi^0 = \theta^0$ and $\phi^1 = \theta^1$. The hope is that  

\textbf{Frozen-features:} Begin by training $\pi_{\theta^0_0}$ and $\pi_{\theta^0_1}$ to produce good CATE estimates for experiment $0$. Assuming $\theta^0_0$ and $\theta^0_1$ have more than $k$ layers, let $\gamma_0$ be the parameters corresponding to the first $k$ layers of $\theta^0_0$. Define $\gamma_1$ analogously. Since we think the features encoded by $\pi_{\gamma_i} (X)$ would make a more informative input than the raw features $X$, we want to use those features as a transformed input space for $\pi_{\theta_0^1}$ and $\pi_{\theta_1^1}$. To wit, set $z_0 =\pi_{\gamma_0}(x)$ and $z_1 = \pi_{\gamma_1}(x)$. Then form the estimates $\pi_{\theta^1_0} (z_0) = \hat{\mu}^1_0$ and $\pi_{\theta_1^1}(z_1) = \hat{\mu}^1_1$. During training of experiment $1$, we only backpropagate through $\theta_0^1$, $\theta_1^1$ and not through the features we borrowed from $\theta_0^0$ and $\theta_1^0$. See Figure \ref{fig:warm-start} and Algorithm \ref{alg:frozen-t} in the appendix.

%Let superscripts represent neural network weights corresponding to a specific network layer. For example, $\theta_0^0$ is the first layer of neural network $\pi_{\theta_0}$. Supposing $\theta_0$ and $\theta_1$ both have $k>2$ layers,
%For experiment 0, let $\theta_0 = \left[\theta_0^0, \theta_0^1, \dots , \theta_0^N \right]$ and $\theta_1 = \left[\theta_1^0, \theta_1^1, \dots , \theta_1^N \right]$ to estimate $\mu^0_0$ and $\mu^0_1$
%For experiment 1, we consider the hybrid network parameterized by $\omega_0 = \left[\theta_0^0, \theta_0^1, \phi_0^0, \phi_0^1, \dots , \phi_0^N \right]$ and likewise for $\omega_1$. Since we think the features learned by $\theta_i$ are good, we choose to not bakpropgate through $\theta_i^j$ during training of $\omega_i$. This is equivalent to transforming the input space of $\pi_{\phi_i}$ from $X$ to $\pi_\gamma(X)$ with $\gamma = \left[ \theta_i^0, \theta_i^1 \right]$. Put another way, $\pi_\phi$ now receives as input not the raw features from $X$ but the features $\pi_\gamma(X)$ that were trained under experiment $0$. See figure ??.

\textbf{Multi-head:} In this setup, all experiments share base layers that are followed by experiment-specific layers. The intuition is that the base layers should learn general features, and the experiment-specific layers should transform those features into estimates of $\mu_j^i$. More concretely, let $\gamma_0$ and $\gamma_1$ be shared base layers. Set $z_0 = \pi_{\gamma_0}(x_0)$ and $z_1 = \pi_{\gamma_1}(x_1)$. The base layers are followed by experiment-specific layers $\phi_0^i$ and $\phi_1^i$. Let $\theta_j^i = \left[\gamma_j, \phi_j^i \right]$. Then $\pi_{\theta_j^i}(x) = \pi_{\phi_j^i} \left(\pi_{\gamma_j}(x) \right) = \pi_{\phi_j^i} (z_j) = \hat{\mu}_j^i$. Training alternates between experiments: each $\theta_0^i$ and $\theta_1^i$ is trained for some small number of iterations, and then the experiment and head being trained are switched. Every head is usually trained several times. See Figure \ref{fig:warm-start} Algorithm \ref{alg:multi-t} in the appendix.


%Let $z_i = \pi_{\gamma_i}(x_i)$. Then we have $\hat{\mu}_0 = \pi_{\omega_0} (X)$. Training alternates between experiments. Each head is trained for some small number of iterations and then the head being trained switches. Each head is usually trained several times. See figure ???.
%$\left[\gamma_i^0, \gamma_i^1, \dots, \gamma_i^k \right]$ $\left[\phi_i^0, \dots, \phi_i^N \right]$

\textbf{Joint training:} All predictions share base layers $\theta$. From these base layers, there are two heads per-experiment $i$: one to predict $\mu_0^i$ and one to predict $\mu_1^i$. Every head and the base features are trained simultaneously by optimizing with respect to the loss function $\mathcal{L} = \sum_i  \| \left(\hat{\mu}_0^i - \mu_0^i \right) \| + \| \left(\hat{\mu}_1^i - \mu_1^i \right) \|$ and minimizing over all weights. This will encourage the base layers to learn generally applicable features and the heads to learn features specific to predicting a single $\mu_j^i$. See Figure Algorithm \ref{alg:joint-training}. 


%The setup is somewhat similar to multi-head with some key differences. First, training no longer alternates. Instead, all experiments are optimized simultaneously. This is accomplished by adding together the loss terms for each experiment $\mathcal{L} = \sum_i  \| \left(\hat{\mu}_0^i - \mu_0^i \right) \| + \| \left(\hat{\mu}_1^i - \mu_1^i \right) \|$ and minimizing over all weights. Second, all $\mu_0^i$ and $\mu_1^i$ are computed directly as in the T-Learner. There are no intermediate calculations as in the X and Y learners. This is in contrast to multi-head, wherein the X and Y learners are permissible and simply receive multiple heads for all required intermediate calculations. Finally, estimates of $\mu_0^i$ and $\mu_1^i$ share only one set of base layers. In multi-head, networks that estimate $\mu_0^i$ and $\mu_1^i$ have distinct base layers.  

%Joint-training is logically equivalent to simultaneously training one T-learners per experiment under the condition that all such T-learners share common base layers which feed into their own independent layers. See algorithm 2 in Appendix 1.1.   

\textbf{SF Reptile transfer for CATE estimators:} Similarly to fine-tuning, we no longer provide each experiment with its own weights. Instead, we use data from all experiments to learn weights $\theta_0$ and $\theta_1$, which are good initializers. By good initializers, we mean that starting from $\theta_0$ and $\theta_1$, one can train neural networks $\pi_{\theta_0}$ and $\pi_{\theta_1}$ to estimate $\mu^i_0$ and $\mu^i_1$ for any arbitrary experiment much faster and with less data than starting from random initializations. To learn these good initializations, we use a transfer learning technique called Reptile. The idea is to perform experiment-specific inner updates $U(\theta)$ and then aggregate them into outer updates of the form $\theta_{\text{new}} = \mathbf{\epsilon} \cdot U(\theta) + (1- \mathbf{\epsilon}) \cdot \theta$. In this paper, we consider a slight variation of Reptile. In standard Reptile, $\epsilon$ is either a scalar or correlated to per-parameter weights furnished via SGD. For our problem, we would like to encourage our network layers to learn at different rates. The hope is that the lower layers can learn more general, slowly-changing features like in the frozen features method, and the higher layers can learn comparatively faster features that more quickly adapt to new tasks after ingesting the stable lower-level features. To accomplish this, we take the path of least resistance and make $\epsilon$ a vector which assigns  a different learning rate to each neural network layer. Because our intuition involves slow and fast weights, we will refer to this modification in this paper as SF Reptile: Slow Fast Reptile. Though this change is seemingly small, we found it boosted performance on our problems. See Figure \ref{fig:joint_training} and Algorithm \ref{alg:sf-reptile-t}.

\textbf{MLRW transfer for CATE estimation:} In this method, there exists one single set of weights $\theta$. There are no experiment-specific weights. Furthermore, we do not use separate networks to estimate $\mu_0$ and $\mu_1$. Instead, $\pi_\theta$ is trained to estimate one $\mu_j^i$ at a time. We train $\theta$ with SF Reptile so that in the future $\pi_\theta$ requires minimal samples to fit $\mu_j^i$ from any experiment. To actually form the CATE estimate, we use a small number of training samples to fit $\pi_\theta$ to $\mu_0^i$ and then a small number of training samples to fit $\pi_\theta$ to $\mu_1^i$. We call $\theta$ \textbf{meta-learned regression weights (MLRW)} because they are meta-learned over many experiments to quickly regress onto any $\mu_j^i$. The full MLRW algorithm is presented as Algorithm \ref{alg:sf-reptile-meta-regression}. 






	
    % Data sources:
% 1.)	http://journals.sagepub.com/doi/abs/10.1177/1532673X16686556 The Generalizability of Social Pressure Effects on Turnout Across High-Salience Electoral Contexts: Field Experimental Evidence From 1.96 Million Citizens in 17 States		
% 2.) https://scholar.harvard.edu/files/todd_rogers/files/social_pressure.pdf Social pressure and voting: A field experiment conducted in a high-salience election
% 3.) https://www.povertyactionlab.org/sites/default/files/publications/4707_The-Comparative-Effectiveness-of-turnout-of-positively-versus-negatively-framed-voter-mob-appeals_March2017.pdf, Study 2 THE COMPARATIVE EFFECTIVENESS ON TURNOUT OF POSITIVELY VERSUS NEGATIVELY FRAMED VOTER MOBILIZATION APPEALS

% TODO:  
% BENCHMARKS: 
% 1.) Reproduce the study suggested by Nie and Wager
% 2.) Design an experiment which shows how powerful it can be to have a big set of people which are neither in the control nor in the treated group
		% It might also be good to not take experiments other people have done, and create experiments
		% which are very closely related to those experiments. (e.g., Demystifying Double Robustness, 2007
		% , section 1.4)
		%
		% CONCLUSIONS we want to draw from these experiments are and should be somehow emphasized 
		% and shown in the experiments are: (We might find that not all of these conclusions are actually true, 
		% but this is what I would expect to be true based on what we have seen so far.)
		% 		a) Neural Networks for heterogeneous treatment effect estimation is a very powerful and
		% 			flexible "base-learner" in other words X-Neural Networks does better than X-RF
		% 		b) Transfer learning for causal inference can significantly improve the prediction performance
		% 		c) (Maybe) Show that Jases suggestion to match known quantities (e.g., average similar to the 
		% 			ATE) will protect us from some forms of over fitting
		%		d) Reptile helps, Multi head helps, but the combination of the two helps the most


\section{Evaluation} \label{section:Simulation}
We evaluate our transfer learning estimators on both real and simulated data.  In our data example, we consider the important problem of voter encouragement. Analyzing a large data set of 1.96 million potential voters, we show how transfer learning across elections and geographic regions can dramatically improve our CATE estimators. This example shows that transfer learning can substantially improve the performance of CATE estimators. \textbf{To the best of our knowledge, this is the first successful demonstration of transfer learning for CATE estimation.} The simulated data has been intentionally chosen to be different in character from our real-world example. In particular, the simulated input space is images and the estimated outcome variable is continuous. %These features make the simulation problem quite challenging.

%\begin{center}
%	\begin{tabular}{ | p{1.7cm} | p{2.3cm} | p{2.3cm} | p{2.3cm} |p{2.3cm} |p{2.3cm} |}
%		\hline
%		Experiment & $x$ & outcome & $\mu_0$ & $\mu_1$ & $\tau$ \\ [0.5ex] 
%		\hline \hline
%		GOTV & a voter profile & \parbox[t]{2.3cm}{The voter's \\ propensity to \\ vote} & \parbox[t]{2.3cm}{The voter's \\ propensity to \\ vote when they \\ do not receive \\ a mailer} & \parbox[t]{2.3cm}{The voter's \\ propensity to \\ vote when they \\ do receive a \\ mailer} & \parbox[t]{2.3cm}{Change in the \\ voter's \\ propensity to \\ vote after \\ receiving a \\ mailer \\ } \\  \hline
%		%MNIST & \parbox[t]{1.7cm}{an MNIST \\ image} & \parbox[t]{2.3cm}{ Price generated \\ for the image. \\ Depends on \\ magnitude of \\ depicted number \\ } & \parbox[t]{2.3cm}{Price function \\ for poorly lit \\ image} & \parbox[t]{2.3cm}{Price function \\ for well lit \\ image} & \parbox[t]{2.3cm}{Difference in \\ price between \\ poorly lit and \\ well lit image} \\ [1ex] 
%		\hline
%	\end{tabular}
%\end{center}

%on a real world experiment and we use two different data sources to improve the estimate of the treatment effect estimates. 

%Performing well on simulations can reveal very useful insights on the relevant effects, but the real important benchmark is how well it performs on real data.




%To emphasize that the conclusions we draw are not specific for voter persuasion studies and the available features in our lead example



\subsection{GOTV Experiment} \label{section:GOTV}

%In a field experiments with 229,461 potential voters, \cite{GerberGreenLarimer} evaluate the effect of a mailer on the voter turnout on the August 2006 primary election in Michigan. The mailer starts with the words ``WHO VOTES IS PUBLIC INFORMATION'' it  contains the voting history for all subjects in the household, a blank field for the upcoming election and the letter promises to sent an updated version after the election. For each individual the study contains seven individual discrete covariates: gender, age, and the voting turnout for the primary elections in 2000, 2002, 2004, and the general election in 2000 and 2002. 

%BCS 05/16: I think we need to make it clear what $\mu_0^i$, $\mu_1^i$, $X^i$, and $\tau^i$ are for this dataset. Also emphasize what the treatment effect we're measuring is and why it matters. Maybe we should even put them in a table or something?

To evaluate transfer learning for CATE estimation on real data, we reanalyze a set of large field experiments with more than 1.96 million potential voters \citep{gerber2017generalizability}. The authors conducted 17 experiments to evaluate the effect of a mailer on voter turnout in the 2014 U.S.\ Midterm Elections. 
The mailer informs the targeted individual whether or not they voted in the past four major elections (2006, 2008, 2010, and 2012), and it compares their voting behavior with that of the people in the same state. The mailer finishes with a reminder that their voting behavior will be monitored. The idea is that social pressure---i.e., the social norm of voting---will encourage people to vote. 
The likelihood of voting increases by about 2.2\% (s.e.=0.001) when given the mailer.


Each of the experiments target a different state. This results in different populations, different ballots, and different electoral environments. In addition to this, the treatment is slightly different in each experiment, as the median voting behavior in each state is different.
However, there are still many similarities across the experiments, so there should be gains from transferring information. %For example, as \cite{kunzel2017meta} point out, potential voters aged 63 to 67 tend to be particularly receptive to this kind of treatment while the effect of such a mailer for younger voters is much smaller.

In this example, the input $X$ is a voter's demographic data including age, past voting turnout in 2006, 2008, 2009, 2010, 2011, 2012, and 2013, marital status, race, and gender. The treatment response function $\hat{\mu}_1(x)$ estimates the voting propensity for a potential voter who receives a mailer encouraging them to vote. The control response function $\hat{\mu}_0$ estimates the voting propensity if that voter did not receive a mailer. 
The CATE $\tau$ is thus the change in the probability of voting when a unit receives a mailer. The complete dataset has this data over 17 different states. Treating each state as a separate experiment, we can perform transfer learning across them.  

\begin{center}
	\begin{tabular}{| p{2.3cm} | p{2.3cm} | p{2.3cm} |p{2.3cm} |p{2.3cm} |}
		\hline
		$x$ & outcome & $\mu_0$ & $\mu_1$ & $\tau$ \\ [0.5ex] 
		\hline \hline
		a voter profile & \parbox[t]{2cm}{The voter's \\ propensity to \\ vote} & \parbox[t]{2cm}{The voter's \\ propensity to \\ vote when they \\ do not receive \\ a mailer} & \parbox[t]{2cm}{The voter's \\ propensity to \\ vote when they \\ do receive a \\ mailer} & \parbox[t]{2cm}{Change in the \\ voter's \\ propensity to \\ vote after \\ receiving a \\ mailer \\ } \\  \hline
		%MNIST & \parbox[t]{1.7cm}{an MNIST \\ image} & \parbox[t]{2cm}{ Price generated \\ for the image. \\ Depends on \\ magnitude of \\ depicted number \\ } & \parbox[t]{2.3cm}{Price function \\ for poorly lit \\ image} & \parbox[t]{2.3cm}{Price function \\ for well lit \\ image} & \parbox[t]{2.3cm}{Difference in \\ price between \\ poorly lit and \\ well lit image} \\ [1ex] 
		\hline
	\end{tabular}
\end{center}

Being able to estimate the treatment effect of sending a mailer is an important problem in elections. We may wish to only treat people whose likelihood of voting would significantly increase when receiving the mailer, to justify the cost for these mailers. Furthermore, we wish to avoid sending mailers to voters who will respond negatively to them. %Applied researchers have observed a backlash from these mailers
This negative response has been previously observed and is therefore feasible and a relevant problem---e.g., some recipients call their Secretary of State's office or local election registrar to complain \citep{mann2010there,michelson2016risk}.  

\subsubsection*{Evaluating CATE estimators on real data}
Evaluating a CATE estimator on real data is difficult since one does not observe the true CATE or the individual treatment effect, $Y_i(1) - Y_i(0)$, for any unit because by definition only one of the two outcomes is observed for any unit. One could use the original features and simulate the outcome features, but this would require us to create a response model. Instead, we estimate the "truth" on the real data using linear models (version 1) or random forests (version 2), and we then draw the data based on these estimates. For a detailed description, we refer to Appendix \ref{section:Appendix-GOTV}. We then ask the question: how do the various methods perform when they have less data than the entire sample?
%JSS: I really don't like the linear thing....kind of sucks.

We evaluate S-NN, T-NN, and Y-NN using our transfer learning methods. We also added a baseline benchmark which does not use any transfer learning for each of the CATE estimators. 
In addition to this, we added the S-RF and T-RF as random forest baselines, as well as the Joint estimator and the MLRW estimator, both of which use transfer learning.
Figure \ref{fig:mainpapergotv21lm1} shows the performance of these estimators when the regression functions were created using a linear model, and Figure \ref{fig:mainpaperGOTV22rf1} shows the same, but the response functions are created using a random forest fitted on the real data. 

In previous work, the non-transfer tree-based estimators such as T-RF and S-RF have achieved state of the art results on this problem \citep{kunzel2017meta}. For CATE estimation, these methods are very competitive baselines \citep{green2012modeling}. Happily for us, even non-transfer neural-network-based learners vastly outperform the prior art. In both examples, non-transfer S-NN, T-NN, and Y-NN learners are better or not much worse than T-RF and S-RF. S-NN and Y-NN perform extremely well in this example. Better still, our transfer learning approaches consistently outperform all classical baselines and non-transfer neural network learners on this benchmark. Positive transfer between experiments is readily apparent. 
%, but we don't believe that either learner is consistently the best, and assert that it depends on the underlying data-generating distribution.  

%getting rid of list 
We find that multi-head, frozen features, and SF are usually the best methods to improve an existing neural network-based CATE estimator.  \textbf{The best estimator is MLRW.} This algorithm consistently converges to a very good solution with very few observations.

%The authors collected age, past voting turnout in 2006, 2008, 2009, 2010, 2011, 2012, 2013, martial status, race dummies, and sex for each unit in the training set. 
%We are using these covariates to target the treatments. 

%%%%%%%%%%%

%%PA.5.17: captions for figures don't have enough info; they should be made self-contained;  e.g. need to include which of the experiments (MNIST or CATE), need to include explanation of what's being plotted, what's observed, what's concluded from these observations



\begin{figure}
	\centering
	\includegraphics[width=.8\linewidth]{figures/main_paper_GOTV_2_1_lm_1}
    \caption{Social Pressure and Voter Turnout, Version 1. Our results far exceed the previous state of the art, which are represented here as S-RF, T-RF, and the baseline method for S-NN and T-NN. Our new methods are Y-NN and the transfer learning methods: warm, frozen, multi-head, joint, SF Reptile, and MLRW.}
    %captions for figures don't have enough info; they should be made self-contained;  e.g. need to include which of the experiments (MNIST or CATE), need to include explanation of what's being plotted, what's observed, what's concluded from these observations}
    	\label{fig:mainpapergotv21lm1}
\end{figure}

\begin{figure}
	\centering
	\includegraphics[width=.8\linewidth]{figures/main_paper_GOTV_2_2_rf_1}
       \caption{Social Pressure and Voter Turnout, Version 2. Our results exceed the previous state of the art results, which are represented here as S-RF, T-RF, and the baseline method for S-NN and T-NN. Our new methods are Y-NN and the transfer learning methods: warm, frozen, multi-head, joint, SF Reptile, and MLRW.}
       %captions for figures don't have enough info; they should be made self-contained;  e.g. need to include which of the experiments (MNIST or CATE), need to include explanation of what's being plotted, what's observed, what's concluded from these observations
    	\label{fig:mainpaperGOTV22rf1}
\end{figure}

%%PA.5.17: how come truth is linear or random forest? is this by design?  I thought the voter thing was on real data? and MNIST definitely isn't linear classification (is it somehow here?); I am rather confused...



\subsection{MNIST Example} \label{section:MNIST}

%\textbf{BCS 05/16: I think we need to make it clear what $\mu_0^i$, $\mu_1^i$, $X^i$, and $\tau^i$ are for this dataset. Also emphasize what the treatment effect we're measuring is and why it's interesting. As it stands, this section really makes zero sense. I don't really understand the interpretation of the image being `well lit' or not.}


In the previous experiment, we observed that the MLRW estimator performed most favorably and transfer learning significantly improved upon the baseline. To confirm that this conclusion is not specific to voter persuasion studies, we consider in this section intentionally a very different type of data.
Recently, \cite{nie2017learning} introduced a simulation study wherein MNIST digits are rotated by some number of degrees $\alpha$; with $\alpha$ furnished via a single data generating process that depends on the value of the depicted digit. They then attempt to do CATE estimation to measure the heterogeneous treatment effect of a digit's label. 

Motivated by this example, we develop a data generating process using MNIST digits wherein transfer learning for CATE estimation is applicable. In our example, the input $X$ is an MNIST image. We have $k$ data generating processes which return different outcomes for each input when given either treatment or control. Thus, under some fixed data generating process, $\mu_0$ represents the outcome when the input image $X$ is given the control, $\mu_1$ represents the outcome when $X$ is given the treatment, and $\tau$ is the difference in outcomes given the placement of $X$ in the treatment or control group. Each data generating process has different response functions ($\mu_0$ and $\mu_1$) and thus different CATEs ($\tau$), but each of these functions only depend on the image label presented in the image $X$. We thus hope that transfer learning could expedite the process of learning features which are indicative of the label. See Appendix \ref{app:simulation} for full details of the data generation process. 
In Figure \ref{fig:MNIST} of Appendix \ref{app:simulation}, we confirm that a transfer learning strategy outperforms its non-transfer learning counterpart, even on image data, and also that MLRW performs well.
%%PA.5.17: I am confused as to what it means to be willing to pay a price for an image; how has that been determined when the data set was generated?


%For our MNIST simulation study (Section \ref{section:MNIST}), we used the MNIST database \citep{lecun1998mnist} which contains labeled handwritten images. We follow here the notation of \cite{nie2017learning}, who introduce a very similar simulation study which is not trying to evaluate transfer learning for CATE estimation, but it also emulates a RCT with the goal to evaluate different CATE estimators. 

%SRK: ! We will have both simulations and we will say that, if the data is maximally different (e.g. simulation 2 MLRW will learn very littel while in very similar situations it will learn a lot. 

%This example is motivated by a recent paper of \cite{nie2017learning}. Here we simulate several randomized controlled trials using image data of handwritten digits as input features \citep{lecun1998mnist}. We then simulate the response functions by rotating the digits. 


















		% TODO:  
		% BENCHMARKS: We will need to benchmark our estimators with the following estimators: X-RF, 
		% S-RF, T-RF, Wager's R learner, S-NN, T-NN, Y-NN, Y*-NN, Only Multi Head, Only Raptile for T
		% EXPERIMENTS:
		% 		1.) Design a new simulation study which is very similar to the actual experiment
		% 		2.) Reproduce the study suggested by Nie and Wager
		%		3.) Design an experiment which shows how powerful it can be to have a big set of people 
		% 			 which are neither in the control nor in the treated group
		% It might also be good to not take experiments other people have done, and create experiments
		% which are very closely related to those experiments. (e.g., Demystifying Double Robustness, 2007
		% , section 1.4)
		%
		% CONCLUSIONS we want to draw from these experiments are and should be somehow emphasized 
		% and shown in the experiments are: (We might find that not all of these conculsions are actually true, 
		% but this is what I would expect to be true based on what we have seen so far.)
		% 		a) Neural Networks for heterogenenous treatment effect estimation is a very powerful and
		% 			flexible "base-learner" in other words X-Neural Networks does better than X-RF
		% 		b) Transferlearning for causal inference can significanly improve the prediciton performance
		% 		c) (Maybe) Show that Jases suggestion to match known quantities (e.g., average similar to the 
		% 			ATE) will protect us from some forms of overfitting
		%		d) Raptlie helps, Multihead helps, but the combination of the two helps the most

	\section{Discussion and Conclusion}

%%PA.5.17: TODO this section? :)

%%PA.5.17: confusing the paper ends with a related work section that's living in the Discussion/Conclusion; make this its own section somewhere?

In this paper, we proposed the problem of transfer learning for CATE estimation. One immediate question the reader may be left with is why we chose the transfer learning techniques we did. We only considered two common types of transfer: 
(1) Basic fine tuning and weights sharing techniques common in the computer vision literature \citep{finetune0,finetune2,finetune1, donahu, koch},
(2) Techniques for learning an initialization that can be quickly optimized \citep{maml,hugo,reptile}. However, many further techniques exist. Yet, transfer learning is an extensively studied and perennial problem \citep{schmid,bengio,thurn2,thrun,taylor,silver}. In \cite{oriol}, the authors attempt to combine feature embeddings that can be utilized with non-parametric methods for transfer. \cite{proto} is an extension of this work that modifies the procedure for sampling examples from the support set during training. \cite{marcin} and related techniques try to meta-learn an optimizer that can more quickly solve new tasks. \cite{progressive} attempts to overcome forgetting during transfer by systematically introducing new network layers with lateral connections to old frozen layers. \cite{yu} uses networks with memory to adapt to new tasks. We invite the reader to review \cite{maml} for an excellent overview of the current transfer learning landscape. Though the majority of the discussed techniques could be extended to CATE estimation, our implementations of \cite{progressive, marcin} proved difficult to tune and consequently learned very little. Furthermore, we were not able to successfully adapt \cite{proto} to the problem of regression. We decided to instead focus our attention on algorithms for obtaining good initializations, which were easy to adapt to our problem and quickly delivered good results without extensive tuning. 

%Though our findings were largely positive, we are still left with several open questions. Most immediately, is it possible to do transfer if the input features are dynamic? This would open up a wealth of datasets wherein the inputs are similar, but not identical. YOUR MOM IS FAT



%SRK18: * We know that there are cases when one learner e.g., the T-learner or the S-learner is better. We see in our experiment that the T-learner can benefit hugely from meta learnerning. The S-learner did already do very well and it was difficult to improve its performance.

%The transfer learning techniques we investigated all involved fairly straightforward weight sharing techniques or techniques for meta learning good initialization weights. However, many other transfer learning algorithms exist. \cite{proto,progressive,schmid,donahu}

	\clearpage
	
	\bibliographystyle{apalike}
	\bibliography{main.bbl}
	
    \clearpage

    \appendix
    \clearpage
    \newcommand{\N}{\mathcal{N}}

\section{Appendix: Simulation Studies and Application}
\FloatBarrier
\begin{figure}[H]
	\centering
	\includegraphics[width=.5\linewidth]{figures/MNIST_20_eval}
    \caption{MNIST task}
    	\label{fig:MNIST}
\end{figure}

\label{app:simulation}
\subsection{MNIST Simulation}
For our MNIST simulation study (Section \ref{section:MNIST}), we used the MNIST database \citep{lecun1998mnist} which contains labeled handwritten images. We follow here the notation of \cite{nie2017learning}, who introduce a very similar simulation study which is not trying to evaluate transfer learning for CATE estimation, but instead emulates a RCT with the goal to evaluate different CATE estimators. 

The MNIST data set contains labeled image data $(X_i, C_i)$, where $X_i$ denotes the raw image of $i$ and $C_i \in \{0, \ldots, 9\}$ denotes its label. 
We create $k$ Data Generating Processes (DGPs), $D_1,~\ldots, ~ D_k$, each of which specifies a distribution of $(Y_i(0),Y_i(1), W_i, X_i)$ and represents different CATE estimation problems. 

In this simulation, we let $W_i = 0$ if the image $X_i$ is placed in the control, and $W_i = 1$ if the image $X_i$ is placed in the treatment. $Y_i(W_i)$ quantifies the the outcome of $X_i$ under $W_i$.

To generate a DGP $D_j$ , we first sample weights in the following way, 
\begin{align*}
m^j(0), ~m^j(1), \ldots, ~m^j(9) &\overset{iid}{\sim} \mbox{Unif}(-3,~ 3),\\ 
t^j(0), ~t^j(  1), \ldots, ~t^j(9) &\overset{iid}{\sim} \mbox{Unif}(-1,~ 1),\\ 
p^j(0), ~p^j(1), \ldots, ~p^j(9) &\overset{iid}{\sim} \mbox{Unif}(0.3,~ 0.7),
\end{align*}
and we define the response functions and the propensity score as
\begin{align*}
\mu^j_0(C_i) & = m^j(C_i) + 3 C_i,\\ 
\mu^j_1(C_i) & = \mu^j_0(C_i) + t^j(C_i),\\
e^j(C_i) & = p^j(C_i).
\end{align*}
To generate $(Y_i(0),Y_i(1), W_i, X_i)$ from $D_j$, we fist sample a $(X_i, C_i)$ from the MNIST data set, and we then generate $Y_i(0),Y_i(1)$, and $W_i$ in the following way:
\begin{align*}
\varepsilon_i &\overset{iid}{\sim} \N(0,1)\\
Y_i(0) &= \mu_0(C_i) + \varepsilon_i\\
Y_i(1) &=  \mu_1(C_i) + \varepsilon_i\\
W_i &\sim  \mbox{Bern}(e(C_i)).
\end{align*}

During training, $X_i, W_i$, and $Y_i$ are made available to the convolutional neural network, which then predicts $\hat{\tau}$ given a test image $X_i$ and a treatment $W_i$. $\tau$ is the difference in the outcome given the difference in treatment and control.

Having access to multiple DGPs can be interpreted as having access to prior experiments done on a similar population of images, allowing us to explore the effects of different transfer learning methods when predicting the effect of a treatment in a new image.




    \subsection{GOTV Data Example and Simulation} \label{section:Appendix-GOTV} \label{sec:simulatedGOTV}
In this section, we describe how the simulations for the GOTV example in the main paper were done and we discuss the results of a much bigger simulation study with 51 experiments which is summarized in Tables \ref{table:LM}, \ref{table:RF}, and \ref{table:RF1}.

\subsubsection{Data Generating Processes for Our Real World Example}
For our data example, we took one of the experiments conducted by \cite{gerber2017generalizability}. The study took place in 2014 in Alaska and 252,576 potential voters were randomly assigned in a control and a treatment group. Subjects in the treatment group were sent a mailer as described in the main text and their voting turnout was recorded. 

To evaluate the performance of different CATE estimators we need to know the true CATEs, which are unknown due to the fundamental problem of causal inference.
To still be able to evaluate CATE estimators researchers usually estimate the potential outcomes using some machine learning method and then generate the data from this estimate. This is to some extend also a simulation, but unlike classical simulation studies it is not up to the researcher to determine the data generating distribution.
The only choice of the researcher lies in the type of estimator she uses to estimate the response functions. To avoid being mislead by artifacts created by a particular method, we used a linear model in \textbf{Real World Data Set 1} and random forests estimator in \textbf{Real World Data Set 2}.

Specifically, we generate for each experiment a true CATE and we simulate new observed outcomes based on the real data in four steps.
\begin{enumerate}
\item We first use the estimator of choice (e.g., a random forests estimator) and train it on the treated units and on the control units separately to get estimates for the response functions, $\mu_0$ and $\mu_1$.
\item Next, we sample $N$ units from the underlying experiment to get the features and the treatment assignment of our samples $(X_i, W_i)_{i=1}^N$.
\item We then generate the true underlying CATE for each unit using $\tau_i = \tau(X_i) = \mu_1(X_i) - \mu_0(X_i)$.
\item Finally we generate the observed outcome by sampling a Bernoulli distributed variable around mean $\mu_i$. 
  \begin{align*}
  	&\Yobs_i \sim \mbox{Bern}(\mu_i),&
    &\mu_i = 
    \begin{cases}
    \mu_0(X_i) ~ ~ \mbox{if} ~ ~ W = 0,\\
    \mu_1(X_i) ~ ~ \mbox{if} ~ ~ W = 1.
    \end{cases}&
  \end{align*}
\end{enumerate}

After this procedure, we have 17 data sets corresponding to the 17 experiments for which we know the true CATE function, which we can now use to evaluate CATE estimators and CATE transfer learners.

\subsubsection{Data Generating Processes for Our Simulation Study}
Simulations motivated by real-world experiments are important to assess whether our methods work well for voter persuasion data sets, but it is important to also consider other settings to evaluate the generalizability of our conclusions. 

To do this, we first specify the control response function, $\mu_0(x) = \E[Y(0)|X = x] \in [0,1]$, 
and the treatment response function, $\mu_1(x) = \E[Y(1)|X = x] \in [0,1]$.

We then use each of the 17 experiments to generate a simulated experiment in the following way: 
\begin{enumerate}
\item We sample $N$ units from the underlying experiment to get the features and the treatment assignment of our samples $(X_i, W_i)_{i=1}^N$.
\item We then generate the true underlying CATE for each unit using $\tau_i = \tau(X_i) = \mu_1(X_i) - \mu_0(X_i)$.
\item Finally we generate the observed outcome by sampling a Bernoulli distributed variable around mean $\mu_i$. 
  \begin{align*}
  	&\Yobs_i \sim \mbox{Bern}(\mu_i),&
    &\mu_i = 
    \begin{cases}
    \mu_0(X_i) ~ ~ \mbox{if} ~ ~ W = 0,\\
    \mu_1(X_i) ~ ~ \mbox{if} ~ ~ W = 1.
    \end{cases}&
  \end{align*}
\end{enumerate}

The experiments range in size from 5,000 units to 400,000 units per experiment and the covariate vector is 11 dimensional and the same as in the main part of the paper.
We will present here three different setup. 

\textbf{Simulation LM} (Table \ref{table:LM}):
We choose here $N$ to be all units in the corresponding experiment.
Sample $\beta^0 = (\beta^0_1, \ldots, \beta^0_d) \overset{iid}{\sim} \N(0,1)$ and
$\beta^1 = (\beta^1_1, \ldots, \beta^1_d) \overset{iid}{\sim} \N(0,1)$ and define,
\begin{align*}
\mu_0(x) &= \mbox{logistic}\left(x \beta^0\right), \\
\mu_1(x) &= \mbox{logistic}\left(x \beta^1\right).
\end{align*}
\\ \textbf{Simulation RF} (Table \ref{table:RF}):
We choose here $N$ to be all units in the corresponding experiment.
\begin{enumerate}
	\item Train a random forests estimator on the real data set and define $\mu_0$ to be the resulting estimator,
	\item Sample a covariate $f$ (e.g., age),
    \item ample a random value in the support of $f$ (e.g., 38),
    \item Sample a shift $s \sim \N(0, 4).$
\end{enumerate}
Now define the potential outcomes as follows:
\begin{align*} 
\mu_0(x) &= \mbox{trained Random Forests algorithm}\\
\mu_1(x) &= \mbox{logistic}\left(\mbox{logit}\left(\mu_0(x) + s * 1_{f\ge v} \right)\right)
\end{align*}
\\ \textbf{Simulation RFt} (Table \ref{table:RF1}):
This experiment is the same as Simulation RF, but use only one percent of the data, $N = \frac{\# units}{100}$.


\subsubsection*{Results of 42 Simulations}
Even though we combine each Simulation setup with 17 experiments, we only report the first 14, because the last three don't add any new insight, but they don't fit well on the page. 
Looking at Tables \ref{table:LM}, \ref{table:RF}, and \ref{table:RF1}, we observe that MLRW is the best performing transfer learner. In fact, for Simulation LM it is the best in 8 out of 17 experiments, in Simulation RF it is the best in 11 out of 17 experiments, and in Simulation RFt it is best in 10 out of 17 experiments. We also notice that in cases, where it is not the best performing estimator, it is usually very close to the best and it does not fail terribly anywhere. 
For the other transfer method, we note that frozen features, multi-head, and SF works very well and consistently improves upon the baseline learners which are not using outside information. Warm Start, however, does not work well and often even leads to worse results than the baseline estimators.  




    \section{Y-learner} \label{section:XvsY}
In this section, we show the favorable behavior of the Y-learner over the X-learner. In order to show this, we implemented the X-learner exactly as it is described in \citep{kunzel2017meta} and the Y-learner as it is described in Algorithm \ref{algo:Ylearner}. 
Figure \ref{fig:compareXY} shows the MSE in proportion to its sample size. We can see that the X--learner is consistently outperformed on all these data sets by the Y-learner. 
We note that all these data sets were intentionally crated to be very similar to the GOTV data set we are interested in studying. Therefore these data sets are not extremely different from each other, and it is possible that the X-NN performs much better on different data sets. 
\FloatBarrier
%\begin{figure}
%	\centering
%	\includegraphics[width=0.5\linewidth]{figures/y_learner}
%    \caption{Y-learner with Neural Networks}
%    \label{fig:ylearner}
%\end{figure}


\begin{figure}
	\centering
	\includegraphics[width=1\linewidth]{figures/compareXY}
    \caption{In this figure, we compare the Y and the X learner on six simulated data sets. A precise description on how the data was created can be found in Section \ref{sec:simulatedGOTV}}
    	\label{fig:compareXY}
\end{figure}

\subsubsection*{The Y-Learner}
Another important advantage of neural networks is that they can be trained jointly. This enables us to adapt well-performing meta-learners to perform even better. 
Specifically, we used the idea of X-NN to propose %%PA.5.17: invent -> propose ?
%%PA.5.17: it's a bit confusing to read about X-NN here, but it's not defined yet (and it's not clear it's defined anywhere actually...); I am assuming X-NN means an X-learner that uses a neural net, but it's nowhere made explicit...
a new CATE estimator, which we call Y-NN.\footnote{Y is chosen as it is the next letter in the alphabet after X. However, this is not a meta-learner because there is no obvious way to extend it to arbitrary base learners, such as RF or BART.}
The X-learner is essentially a two step procedure.
In the first stage, the outcome functions, $\hat \mu_0$ and $\hat \mu_1$, are estimated and the individual treatment effects are imputed: 
\begin{align*}
&D_i^1 \define Y(1) - \hat \mu_0(X_i)& &\mbox{and} &D_i^0 \define \hat \mu_1(X_i) - Y_i(0).& 
\end{align*}
In the second stage, estimators for the CATE are derived by regressing the features $X$ on the imputed treatment effects. \cite{kunzel2017meta} provides details. In the X-learner, the estimators of the first stage are held fixed and are not updated in the second stage. This is necessary since, unlike neural networks, many machine learning algorithms, such as RF and BART, cannot be updated in a meaningful way once they have been trained. For neural networks and similar gradient optimization-based algorithms, it is possible to jointly update the estimators in the first and the second stage. 
%%PA.5.17: how is this used later to get a CATE estimate? [nowhere explained I think...]  

This is exactly the motivation of the Y-learner. Instead of first deriving an estimator for the control response functions and then an estimator for the CATE function, these functions are optimized jointly. The pseudo-code in Algorithm \ref{algo:Ylearner} shows how these two stages are updated simultaneously. 
In Figure \ref{fig:compareXY}, we compare Y-NN with X-NN, and we find that Y-NN outperforms X-NN for our data sets.

	\clearpage
    \newcommand{\X}{\mathcal{X}}
\newcommand{\I}{\mathcal{I}}
\newcommand{\J}{\mathcal{J}}
%\newcommand{\E}{\mathbb{E}}
\renewcommand{\P}{\mathcal{P}}
%\newcommand{\Yobs}{Y}
\newcommand{\Yobst}{Y^{obs,t}}
\newcommand{\Yobsc}{Y^{obs,c}}
\newcommand{\var}{\mbox{Var}}
\newcommand{\sd}{\mbox{Sd}}
\newcommand{\cov}{\mbox{Cov}}
%\newcommand{\R}{\mathbb{R}}
\newcommand{\Dt}{\tilde D}
\newcommand{\taun}{\hat{\tau}^{mn}}
\newcommand{\theoremfont}{\itshape}
\renewcommand{\descriptionlabel}[1]{\hspace{\labelsep}\textbf{\descriptionlabelfont #1}}
\newcommand{\D}{\mathcal{D}}
\newcommand{\trace}{\mbox{trace}}
\newcommand{\IMSE}{\mbox{IMSE}}
\newcommand{\EMSE}{\mbox{EMSE}}
\newcommand{\MSE}{\mbox{MSE}}
\newcommand{\BIAS}{\mbox{BIAS}}
%\newcommand{\Pcal}{\mathcal{P}}
\newcommand{\mae}{\gamma_{\scriptsize\mbox{max}}}
\newcommand{\mie}{\gamma_{\scriptsize\mbox{min}}}
\newcommand{\mis}{s_{\scriptsize\mbox{min}}}
\newcommand{\ED}{\mathbb{E}_{\mathcal{D}}}
\newcommand{\Ex}{\mathbb{E}_{x}}
\newcommand{\Pobs}{\mathcal{P}_{\mbox{\tiny observed}}}
%\newcommand{\define}{:=}
\newcommand{\deltat}{\tilde \delta}
%\newcommand{\define}{\overset{\mbox{def}}{=}}
%\newcommand{\emin}{e_{\min}}
%\newcommand{\emax}{e_{\max}}
\newcommand{\Sb}{{\bar S}}
\section{Pseudo Code for CATE Estimators}
\label{app:learners}

In this section, we will present pseudo code for the CATE estimators in this paper. We present code for the meta learning algorithms in Section \ref{app:transfer.algos}.
We denote by $Y^0$ and $Y^1$ the observed outcomes for the control and the treated group. For example, $Y^1_i$ is the observed outcome of the $i$th unit in the treated group. $X^0$ and $X^1$ are the features of the control and treated units, and hence, $X^1_i$  corresponds to the feature vector of the $i$th unit in the treated group. $M_k(Y\sim X)$ is  the notation for a regression estimator, which estimates $x \mapsto \E[Y|X=x]$. It can be any regression/machine learning estimator, but in this paper we only choose it to be a neural network or random forest.

% cateestimate T-learner    	
\begin{algorithm}[H] 
	\begin{algorithmic}[1]
		\Procedure{T--Learner}{$X, \Yobs, W$}
		
		\State $\hat \mu_0 	= M_0(Y^0 \sim X^0)$
		\State $\hat \mu_1	 = M_1(Y^1 \sim X^1)$
		
		\item[]
		\State $\hat \tau(x) = \hat \mu_1(x) - \hat \mu_0(x)$ 
		\EndProcedure
	\end{algorithmic}
	\caption{T-learner}
	
	\label{algo:Tlearner}
		\algcomment{$M_0$ and $M_1$ are here some, possibly different machine learning/regression algorithms.}	
\end{algorithm}

% cateestimate S-learner
\begin{algorithm}[H] 
	\begin{algorithmic}[1]
		\Procedure{S--Learner}{$X, \Yobs, W$}
		
		\State $\hat \mu 	= M(\Yobs \sim (X, W))$
%				\Comment{Estimate $\mu_1$ and $\mu_0$}:
						
		\State $\hat \tau(x) = \hat \mu(x,1) - \hat \mu(x,0)$ 
%				\Comment{Difference is the CATE estimator}:
		\EndProcedure
	\end{algorithmic}
	\caption{S-learner}
	
	\label{algo:Slearner}
	\algcomment{$M(\Yobs \sim (X, W))$ is the notation for estimating $(x,w) \mapsto \E[Y|X=x, W=w]$ while treating $W$ as a 0,1--valued feature.}
\end{algorithm}

\begin{algorithm}[H]
	\begin{algorithmic}[1]	
		\Procedure{X--Learner}{$X, \Yobs, W, g$}
		
		
		\item[]
		\State $\hat \mu_0	= M_1(Y^0 \sim X^0)$ 
				\Comment{Estimate response function}
		\State $\hat \mu_1	 = M_2(Y^1 \sim X^1)$
		
		\item[]
		\State $\Dt^1_{i} =         Y^1_{i}                         -     \hat \mu_0(X^1_{i})$
				\Comment{Compute imputed treatment effects}
		\State $\Dt^0_{i} =     \hat \mu_1(X^0_{i})     -     Y^0_{i}$
		
		\item[]
		\State $\hat \tau_1      = M_3(\Dt^{1}      \sim     X^1)$
				\Comment{Estimate CATE for treated and control}
		\State $\hat \tau_0     = M_4(\Dt^0     \sim     X^0)$
		
		\item[]
		\State $\hat \tau(x) = g(x) \hat \tau_0(x) + (1- g(x)) \hat \tau_1(x)$ 
				\Comment{Average the estimates}
		
		\EndProcedure
	\end{algorithmic}
	\caption[test]{X--learner}

	\label{aglo:Xlearner}
	
	\algcomment{$g(x) \in [0,1]$ is a weighing function which is chosen to minimize the variance of $\hat \tau(x)$. It is sometimes possible to estimate $\cov(\tau_0(x) , \tau_1(x) )$, and compute the best $g$ based on this estimate. However, we have made good experiences by choosing $g$ to be an estimate of the propensity score, but also choosing it to be constant and equal to the ratio of treated units usually leads to a good estimator of the CATE. }	
\end{algorithm}


\begin{algorithm}[t]
  \begin{algorithmic}[1]
  \If{$W_i == 0$}
  \State Update the network $\pi_{\theta_0}$ to predict $Y_i^{obs}$
  \State Update the network $\pi_{\theta_1}$ to predict $Y_i^{obs} + \pi_\tau(X_i)$
  \State Update the network $\pi_\tau$ to predict $\pi_{\theta_1}(X_i) - Y_i^{obs}$
  \EndIf
  \If{$W_i == 1$}
  \State Update the network $\pi_{\theta_0}$ to predict $Y_i^{obs} - \pi_\tau(X_i)$
  \State Update the network $\pi_{\theta_1}$ to predict $Y_i^{obs}$
  \State Update the network $\pi_\tau$ to predict $Y_i^{obs} - \pi_{\theta_0}(X_i)$
  \EndIf
  \end{algorithmic}
  \caption{Y-Learner Pseudo Code}
  \label{algo:Ylearner}
  \label{ylearner}
      \algcomment{This process describes training the Y-Learner for one step given a data point $(Y_i^{obs}, X_i, W_i)$}
\end{algorithm}

	\section{Explicit Transfer Learning Algorithms for CATE Estimation}
\label{app:transfer.algos}

%In section 3.2 we outlined several possible transfer learning algorithms for CATE estimation. However, things were mostly presented at a high level so as to remain agnostic to one's choice of underlying CATE estimator. In this section, we provide explicit algorithms for T and Y transfer learners, and explicitly outline the Joint-Training method. 


\subsection{MLRW Transfer for CATE Estimation}

%99.9 percent sure this will have to live in the appendix. It's a huge space hog. 


\begin{algorithm}[H]
	\caption{MLRW Transfer for Cate Estimation.}
	\label{alg:sf-reptile-meta-regression}
	\begin{algorithmic}[1]
		\State Let $\mu_0^{(i)}$ and $\mu_1^{(i)}$ be the outcome under treatment and control for experiment $i$. 
		\State Let numexps be the number of experiments. 
		\State Let $\pi_\theta$ be an $N$ layer neural network parameterized by $\theta = \left[\theta_0, \dots, \theta_{N} \right]$.
		\State Let $\epsilon = \left[\epsilon_0, \dots, \epsilon_{N} \right]$ be a vector, where $N$ is the number of layers in $\pi_\theta$. 
		\State Let outeriters be the total number of training iterations. 
		\State Let inneriters be the number of inner loop training iterations.
		\For{oiter $<$ outeriters}
		\For{i $<$ numexps}
		\State Sample $X_0$ and $X_1$: control and treatment units from experiment $i$
		\For{$j = [0, 1]$} \Comment{j iterating over treatment and control}
		\State Let $U_0(\theta) = \theta$
		\For{k < inneriters}
		\State $\mathcal{L} = \| \pi_{U_k(\theta)} (X_j) - \mu_j(X_j) \|$
		\State Compute $\nabla_\theta \mathcal{L}$. 
		\State Use ADAM with $\nabla_\theta \mathcal{L}$ to obtain $U_{k+1}(\theta)$. 
		\State $U_k (\theta) = U_{k+1}(\theta)$
		\EndFor
		\For{ $p < N$ }
		\State $\theta_p = \epsilon_p \cdot U_k(\theta_p) + (1 - \epsilon_p) \cdot \theta_p$.  
		\EndFor
		\EndFor
		\EndFor
		\EndFor
		\State To Evaluate CATE estimate, \textbf{do}
		\State $C = []$
		\For{i $<$ numexps}
		\State Sample $X_0$ and $X_1$: control and treatment units from experiment $i$
		\State Sample $X$: test units from experiment $i$. 
		\For{$j$ = [0, 1]} \Comment{j iterating over treatment and control}
		\For{$k <$ innteriters}
		\State $\mathcal{L} = \| \pi_{U_k(\theta)} (X_j) - \mu_j(X_j) \|$
		\State Compute $\nabla_\theta \mathcal{L}$. 
		\State Use ADAM with $\nabla_\theta \mathcal{L}$ to obtain $U_{k+1}(\theta)$. 
		\State $U_k (\theta) = U_{k+1}(\theta)$
		\EndFor
		\State $\hat{\mu}_j = \pi_{U_k(\theta)} (X)$
		\EndFor
		\State  $\hat{\tau}_i = \hat{\mu}_0 - \hat{\mu_1}$
		\State C.append$(\hat{\tau}_i)$
		\EndFor
		\State {\bf return} $C$
	\end{algorithmic}
\end{algorithm}

\clearpage
\subsection{Joint Training}
\FloatBarrier
\begin{figure}
	\begin{minipage}[c]{0.45\textwidth}
	\includegraphics[width=0.8\linewidth]{figures/joint_training}
	\end{minipage}\hfill
	\begin{minipage}[c]{0.4\textwidth}
		\caption{Joint Training - Unlike the Multi-head method which differentiates base layers for treatment and control, the Joint Training method has all observations and experiments (regardless of treatment and control) share the same base network, which extracts general low level features from the data.} \label{fig:joint_training}
	\end{minipage}
\end{figure}

\begin{algorithm}[h]
	\begin{algorithmic}[1]
		\State Let $\mu_0^{(i)}$ and $\mu_1^{(i)}$ be the outcome under control and treatment for experiment $i$. 
		\State Let numexps be the number of experiments. 
		\State Let $\pi_\rho$ be a generic expression for a neural network parameterized by $\rho$. 
		\State Let $\theta$ be base neural network layers shared by all experiments. 
		\State Let $\phi_0^{(i)}$  be neural network layers predicting $\mu_0^{(i)}$ in experiment $i$.
		\State Let $\phi_1^{(i)}$  be neural network layers predicting $\mu_1^{(i)}$ in experiment $i$.
		\State Let $\omega_0^{(i)} = \left[\theta, \phi_0^{(i)} \right]$ be the full prediction network for $\mu_0$ in experiment $i$.
		\State Let $\omega_1^{(i)} = \left[\theta, \phi_1^{(i)} \right]$ be the full prediction network for $\mu_1$ in experiment $i$.
		\State Let $\Omega = \bigcup_{j=0}^{\text{1}} \bigcup_{i=1}^{\text{numexps}} \omega_j^{(i)}$ be all trainable parameters. 
		\State Let numiters be the total number of training iterations
		\For{iter $<$ numiters}
		\State $\mathcal{L}$ = 0
		\For{i $<$ numexps}
		\State Sample $X_0$ and $X_1$: control and treatment units from experiment $i$
		\For{$j$ = [0, 1]} \Comment{j iterating over treatment and control}
		\State $\mathcal{L}_j^{(i)} = \| \pi_{\omega_j^{(i)}} (X_j) - \mu_j(X_j) \|$
		\State $\mathcal{L} = \mathcal{L} + \mathcal{L}_j^{(i)}$
		\EndFor
		\EndFor
		\State Compute $\nabla_\Omega \mathcal{L} = \frac{\partial \mathcal{L}}{\partial \Omega} = \sum_i \sum_j \frac{\partial \mathcal{L}_j^{(i)}}{\partial \omega_j^{(i)}}$
		\State Apply ADAM with gradients given by $\nabla_\Omega \mathcal{L}$. 
		\For{i $<$ numexps}
		\State Sample $X$: test units from experiment $i$
		\EndFor
		\EndFor
		\State $\hat{\mu}_0 = \pi_{\omega_0^{(i)}} (X)$
		\State $\hat{\mu}_1 = \pi_{\omega_1^{(i)}}(X)$
		\State {\bf return} CATE estimate $\hat{\tau} = \hat{\mu}_1 - \hat{\mu}_0$
	\end{algorithmic}
	\caption{Joint Training}
\label{alg:joint-training}
\end{algorithm}

\FloatBarrier

\clearpage
\subsection{T-learner Transfer CATE Estimators}
Here, we present full pseudo code for the algorithms from Section 3 using the T-learner as a base learner. All of these algorithms can be extended to other learners including $S, R, X,$ and $Y$. See the released code for implementations. 

\begin{algorithm}[H]
	\caption{Vanilla T-learner (also referred to as Baseline T-learner)}
	\label{alg:vanilla-t}
	\begin{algorithmic}[1]
		\State Let $\mu_0$ and $\mu_1$ be the outcome under treatment and control. 
		\State Let $X$ be the experimental data. Let $X_t$ be the test data. 
		\State Let $\pi_{\theta_0}$ and $\pi_{\theta_1}$ be a neural networks parameterized by $\theta_0$ and $\theta_1$. 
		\State Let $\theta = \theta_0 \cup \theta_1$. 
		\State Let numiters be the total number of training iterations.
		\State Let batchsize be the number of units sampled. We use 64. 
		\For{i $<$ numiters}
		\State Sample $X_0$ and $X_1$: control and treatment units. Sample batchsize units.  
		\State $\mathcal{L}_0= \| \pi_\theta (X_0) - \mu_0(X_0) \| $
		\State $\mathcal{L}_1 = \| \pi_\theta (X_1) - \mu_1(X_1) \| $
		\State $\mathcal{L} = \mathcal{L}_0 + \mathcal{L}_1$
		\State Compute $\nabla_{\theta} \mathcal{L} = \frac{\partial \mathcal{L}}{\partial \theta}$.
		\State Apply ADAM with gradients given by $\nabla_{\theta} \mathcal{L}$. 
		\EndFor
		\State $\hat{\mu}_0 = \pi_{\theta_0} (X_t)$
		\State $\hat{\mu}_1 = \pi_{\theta_1} (X_t)$
		\State {\bf return} CATE estimate $\hat{\tau} = \hat{\mu}_1 - \hat{\mu}_0$
	\end{algorithmic}
\end{algorithm}


\begin{algorithm}[H]
	\caption{Warm Start T-learner}
	\label{alg:warm-t}
	\begin{algorithmic}[1]
		\State Let $\mu^{i}_0$ and $\mu^{i}_1$ be the outcome under treatment and control for experiment $i$. 
		\State Let $X^i$ be the data for experiment $i$. Let $X^i_t$ be the test data for experiment $i$. 
		\State Let $\pi_{\theta_0}$ and $\pi_{\theta_1}$ be a neural networks parameterized by $\theta_0$ and $\theta_1$. 
		\State Let $\theta = \theta_0 \cup \theta_1$. 
		\State Let numiters be the total number of training iterations.
		\State Let batchsize be the number of units sampled. We use 64. 
		\For{i $<$ numiters}
		\State Sample $X^0_0$ and $X^0_1$: control and treatment units for experiment $0$. Sample batchsize units.  
		\State $\mathcal{L}_0= \| \pi_{\theta_0} (X^0_0) - \mu_0(X^0_0) \| $
		\State $\mathcal{L}_1 = \| \pi_{\theta_1} (X^0_1) - \mu_1(X^0_1) \| $
		\State $\mathcal{L} = \mathcal{L}_0 + \mathcal{L}_1$
		\State Compute $\nabla_{\theta} \mathcal{L} = \frac{\partial \mathcal{L}}{\partial \theta}$.
		\State Apply ADAM with gradients given by $\nabla_{\theta} \mathcal{L}$. 
		\EndFor
		\For{i $<$ numiters}
		\State Sample $X^1_0$ and $X^1_1$: control and treatment units for experiment $1$. Sample batchsize units.  
		\State $\mathcal{L}_0= \| \pi_{\theta_0} (X^1_0) - \mu_0(X^1_0) \| $
		\State $\mathcal{L}_1 = \| \pi_{\theta_1} (X^1_1) - \mu_1(X^1_1) \| $
		\State $\mathcal{L} = \mathcal{L}_0 + \mathcal{L}_1$
		\State Compute $\nabla_{\theta} \mathcal{L} = \frac{\partial \mathcal{L}}{\partial \theta}$.
		\State Apply ADAM with gradients given by $\nabla_{\theta} \mathcal{L}$. 
		\EndFor
		\State $\hat{\mu}_0 = \pi_{\theta_0} (X^1_t)$
		\State $\hat{\mu}_1 = \pi_{\theta_1} (X^1_t)$
		\State {\bf return} CATE estimate $\hat{\tau} = \hat{\mu}_1 - \hat{\mu}_0$
	\end{algorithmic}
\end{algorithm}

\newpage

\begin{algorithm}[H]
	\caption{Frozen Features T-learner}
	\label{alg:frozen-t}
	\begin{algorithmic}[1]
		\State Let $\mu^{i}_0$ and $\mu^{i}_1$ be the outcome under treatment and control for experiment $i$. 
		\State Let $X^i$ be the data for experiment $i$. Let $X^i_t$ be the test data for experiment $i$. 
		\State Let $\pi_\rho$ be a generic expression for a neural network parameterized by $\rho$. 
		\State Let $\theta_0^0, \theta_0^1, \theta_1^0, \theta_1^1$ be neural network parameters. The subscript indicates the outcome that $\theta$ is associated with predicting (0 for control and 1 for treatment) and the superscript indexes the experiment.
		\State Let $\gamma_0$ be the first $k$ layers of $\pi_{\theta^0_0}$. Define $\gamma_1$ analogously. 
		\State Let $\theta^i = \theta^i_0 \cup \theta^i_1$. 
		\State Let numiters be the total number of training iterations.
		\State Let batchsize be the number of units sampled. We use 64. 
		\For{i $<$ numiters}
		\State Sample $X^0_0$ and $X^0_1$: control and treatment units for experiment $0$. Sample batchsize units.  
		\State $\mathcal{L}_0= \| \pi_{\theta^0_0} (X^0_0) - \mu_0(X^0_0) \| $
		\State $\mathcal{L}_1 = \| \pi_{\theta^0_1} (X^0_1) - \mu_1(X^0_1) \| $
		\State $\mathcal{L} = \mathcal{L}_0 + \mathcal{L}_1$
		\State Compute $\nabla_{\theta} \mathcal{L} = \frac{\partial \mathcal{L}}{\partial \theta}$.
		\State Apply ADAM with gradients given by $\nabla_{\theta^0} \mathcal{L}$. 
		\EndFor
		\For{i $<$ numiters}
		\State Sample $X^1_0$ and $X^1_1$: control and treatment units for experiment $1$. Sample batchsize units. 
		\State Compute $Z^1_0 = \pi_\gamma(X^1_0)$ and $Z^1_1 = \pi_\gamma(X^1_1)$
		\State $\mathcal{L}_0= \| \pi_{\theta_0^1} (Z^1_0) - \mu_0(Z^1_0) \| $
		\State $\mathcal{L}_1 = \| \pi_{\theta_1^1} (Z^1_1) - \mu_1(Z^1_1) \| $
		\State $\mathcal{L} = \mathcal{L}_0 + \mathcal{L}_1$
		\State Compute $\nabla_{\theta^1} \mathcal{L} = \frac{\partial \mathcal{L}}{\partial \theta^1}$. Do not compute gradients with respect to $\theta^0$ parameters. 
		\State Apply ADAM with gradients given by $\nabla_{\theta^1} \mathcal{L}$. 
		\EndFor
		\State Compute $Z^1_t = \pi_\gamma(X^1_t)$. 
		\State $\hat{\mu_0} = \pi_{\theta^1_0} (Z^1_t)$
		\State $\hat{\mu_1} = \pi_{\theta^1_1} (Z^1_t)$
		\State {\bf return} CATE estimate $\hat{\tau} = \hat{\mu_1} - \hat{\mu_0}$
	\end{algorithmic}
\end{algorithm}

\newpage
\newpage 
\begin{algorithm}[H]
	\caption{Multi-Head T-learner}
	\label{alg:multi-t}
	\begin{algorithmic}[1]
		\State Let $\mu^{i}_0$ and $\mu^{i}_1$ be the outcome under treatment and control for experiment $i$. 
		\State Let $X^i$ be the data for experiment $i$. Let $X^i_t$ be the test data for experiment $i$. 
		\State Let $\pi_\rho$ be a generic expression for a neural network parameterized by $\rho$. 
		\State Let $\theta_0$ be base neural network layers shared by all experiments for predicting outcomes under control.
		\State Let $\theta_1$ be base neural network layers shared by all experiments for predicting outcomes under treatment.
		\State Let $\phi_0^{(i)}$  be neural network layers receiving $\pi_{\theta_0}(x^i_0)$ as input and predicting $\mu_0^{(i)}(x^i_0)$ in experiment $i$.
		\State Let $\phi_1^{(i)}$  be neural network layers receiving $\pi_{\theta_1}(x^i_1)$ as input and predicting $\mu_1^{(i)}(x^i_1)$ in experiment $i$.
		\State Let $\omega_0^{(i)} = \left[\theta, \phi_0^{(i)} \right]$ be all trainable parameters used to predict $\mu_0^i$.
		\State Let $\omega_1^{(i)} = \left[\theta, \phi_1^{(i)} \right]$ be all trainable parameters used to predict $\mu_1^i$.
		\State Let $\Omega^i = \omega_0^{(i)} \cup \omega_1^{(i)}$. 
		\State Let numiters be the total number of training iterations.
		\State Let batchsize be the number of units sampled. We use 64. 
		\State Let numexps be the number of experiments. 
		\For{i $<$ numiters}
		\For{j $<$ numexps}
		\State Sample $X^j_0$ and $X^j_1$: control and treatment units for experiment $j$. Sample batchsize units.  
		\State Compute $Z_0^j = \pi_{\theta_0}(X^j_0)$ and $Z_1^j = \pi_{\theta_1}(X^j_1)$
		\State Compute $\hat{\mu}_0^j = \pi_{\phi_0^j} (z_0^j)$ and $\hat{\mu}_1^j = \pi_{\phi_1^j} (z_1^j)$
		\State $\mathcal{L}_0= \| \hat{\mu}_0^j - \mu_0^j(X^j_0) \| $
		\State $\mathcal{L}_1 = \| \hat{\mu}_1^j - \mu_1^j(X^j_1) \| $
		\State $\mathcal{L} = \mathcal{L}_0 + \mathcal{L}_1$
		\State Compute $\nabla_{\Omega^i} \mathcal{L} = \frac{\partial \mathcal{L}}{\partial \Omega^i}$.
		\State Apply ADAM with gradients given by $\nabla_{\theta} \mathcal{L}$. 
		\EndFor
		\EndFor
		\State Let $C = []$
		\For{j $<$ numexps}
		\State Compute $Z_0^j = \pi_{\theta_0}(X^j_t)$ and $Z_1^j = \pi_{\theta_1}(X^j_t)$
		\State Compute $\hat{\mu}_0^j = \pi_{\phi_0^j} (z_0^j)$ and $\hat{\mu}_1^j = \pi_{\phi_1^j} (z_1^j)$
		\State Estimate CATE $\hat{\tau} = \hat{\mu}_1^j - \hat{\mu}_0^j$. 
		\State $C.$append$(\hat{\tau})$
		\EndFor
		\State {\bf return} $C$ 
	\end{algorithmic}
\end{algorithm}

\newpage

\begin{algorithm}[H]
	\caption{SF Reptile T-learner}
	\label{alg:sf-reptile-t}
	\begin{algorithmic}[1]
		\State Let $\mu_0^i$ and $\mu_1^i$ be the outcome under treatment and control for experiment $i$. 
		\State Let $X^i$ be the data for experiment $i$. Let $X^i_t$ be the test data for experiment $i$. 
		\State Let $\pi_{\theta_0}$ and $\pi_{\theta_1}$ be a neural networks parameterized by $\theta_0$ and $\theta_1$. 
		\State Let $\theta = \theta_0 \cup \theta_1$. 
		\State Let $\epsilon = \left[\epsilon_0, \dots, \epsilon_{N} \right]$ be a vector, where $N$ is the number of layers in $\pi_{\theta_i}$. 
		\State Let numouteriters be the total number of outer training iterations.
		\State Let numinneriters be the total number of inner training iterations. 
		\State Let numexps be the number of experiments. 
		\State Let batchsize be the number of units sampled. We use 64. 
		\For{iouter $<$ numouteriters}
		\For{i $<$ numexps}
		\State $U_0(\theta_0) = \theta_0$ 
		\State $U_0(\theta_1) = \theta_1$. 
		\For{k$<$ numinneriters}
		\State Sample $X_0^i$ and $X_1^i$: control and treatment units. Sample batchsize units.  
		\State $\mathcal{L}_0= \| \pi_{U_k(\theta_0)} (X_0^i) - \mu_0(X_0^i) \| $
		\State $\mathcal{L}_1 = \| \pi_{U_k(\theta_1)} (X_1^i) - \mu_1(X_1^i) \| $
		\State $\mathcal{L} = \mathcal{L}_0 + \mathcal{L}_1$
		\State Compute $\nabla_{\theta} \mathcal{L} = \frac{\partial \mathcal{L}}{\partial \theta}$.
		\State Use ADAM with gradients given by $\nabla_{\theta} \mathcal{L}$ to obtain $U_{k+1}(\theta_0)$ and $U_{k+1}(\theta_1)$. 
		\State Set $U_{k}(\theta_0)=U_{k+1}(\theta_0)$ and $U_{k}(\theta_1) = U_{k+1}(\theta_1)$
		\EndFor
		\For{ $p < N$ }
		\State $\theta_p = \epsilon_p \cdot U_k(\theta_p) + (1 - \epsilon_p) \cdot \theta_p$.  
		\EndFor
		\EndFor
		\EndFor
		\State To Evaluate CATE estimate, \textbf{do}
		\State $C = []$. 
		\For{i $<$ numexps}
		\State $U_0(\theta_0) = \theta_0$ 
		\State $U_0(\theta_1) = \theta_1$. 
		\For{k$<$ numinneriters}
		\State Sample $X_0^i$ and $X_1^i$: control and treatment units. Sample batchsize units.  
		\State $\mathcal{L}_0= \| \pi_{U_k(\theta_0)} (X_0^i) - \mu_0(X_0^i) \| $
		\State $\mathcal{L}_1 = \| \pi_{U_k(\theta_1)} (X_1^i) - \mu_1(X_1^i) \| $
		\State $\mathcal{L} = \mathcal{L}_0 + \mathcal{L}_1$
		\State Compute $\nabla_{\theta} \mathcal{L} = \frac{\partial \mathcal{L}}{\partial \theta}$.
		\State Use ADAM with gradients given by $\nabla_{\theta} \mathcal{L}$ to obtain $U_{k+1}(\theta_0)$ and $U_{k+1}(\theta_1)$. 
		\State Set $U_{k}(\theta_0)=U_{k+1}(\theta_0)$ and $U_{k}(\theta_1) = U_{k+1}(\theta_1)$
		\EndFor
		\State $\hat{\mu}_0^i = \pi_{U_k(\theta_0)}(X_0^i)$
		\State $\hat{\mu}_1^i = \pi_{U_k(\theta_1)}(X_1^i)$
		\State $\hat{\tau}^i =  \hat{\mu}_1^i - \hat{\mu}_0^i$
		\State $C.$append$(\hat{\tau}^i)$. 
		\EndFor
		\State {\bf return} $C$. 
	\end{algorithmic}
\end{algorithm}


    \clearpage
	% latex table generated in R 3.4.4 by xtable 1.8-2 package
% Thu May 17 18:43:51 2018
\begin{sidewaystable}[ht]
\begin{flushleft}
\begin{tabular}{rrrrrrrrrrrrrrrr}
 {} & {\textbf{Method}} & {\textbf{LM-1}} & {\textbf{LM-2}} & {\textbf{LM-3}} & {\textbf{LM-4}} & {\textbf{LM-5}} & {\textbf{LM-6}} & {\textbf{LM-7}} & {\textbf{LM-8}} & {\textbf{LM-9}} & {\textbf{LM-10}} & {\textbf{LM-11}} & {\textbf{LM-12}} & {\textbf{LM-13}} & {\textbf{LM-14}} \\ 
  \hline
{\textbf{t-lm}} &  & 15.95 & 7.82 & 20.14 & 6.62 & 46.15 & 17.24 & 9.88 & 10.65 & 44.77 & 7.63 & 8.43 & 9 & 11.21 & 20.9 \\ 
   \hline
{\textbf{s-rf}} &  & 19.13 & 13.18 & 7.62 & 13.27 & 15.66 & 15.58 & 11.44 & 16.3 & 11.34 & 16.53 & 12.57 & 14.11 & 13.49 & 18.56 \\ 
   \hline
{\textbf{t-rf}} &  & 20.04 & 13.56 & 7.79 & 13.62 & 16 & 15.99 & 11.76 & 16.65 & 11.54 & 17.11 & 13.66 & 14.38 & 13.67 & 18.96 \\ 
   \hline
{\textbf{R-NN}} &  & 18.95 & 5.19 & 9.33 & 28.98 & 3.04 & 10.03 & 4.88 & 16.34 & 7.91 & 14.56 & 19.23 & 3.5 & 10.4 & 15.53 \\ 
   \hline
{\textbf{S-NN}} & frozen & 6.74 & 6.17 & 2.75 & 5.76 & 5.68 & 6.27 & 4.23 & 7.17 & 4.15 & 6.77 & 4.88 & 5.9 & 6.63 & 9.87 \\ 
   & multi head & 3.3 & 3.05 & 2.34 & 3.83 & 3.6 & 4.89 & 4.56 & 8.47 & 5.87 & 5.22 & 5.43 & 6.77 & 5.15 & 5.58 \\ 
   & SF & 4.85 & 38.65 & 7.54 & 55.92 & 11.36 & 3.98 & 50.76 & 7.74 & 5.72 & 7.73 & 8.64 & 31.06 & 30.81 & 18.08 \\ 
   & baseline & 6.84 & 5.29 & 3.77 & 6.86 & 5.94 & 7.19 & 4.6 & 8.08 & 5.61 & 7.12 & 4.34 & 6.05 & 8.59 & 10.75 \\ 
   & warm & 7.29 & 6.44 & 4.08 & 7.25 & 6.74 & 7.17 & 6.03 & 8.8 & 5.63 & 7.6 & 5.82 & 6.53 & 8.55 & 12.02 \\ 
   \hline
{\textbf{T-NN}} & frozen & 6.53 & 6.73 & 4.49 & 6.35 & 7.23 & 6.61 & 5.99 & 7.59 & 5.79 & 6.79 & 5.99 & 6.58 & 7.38 & 9.39 \\ 
   & multi head & 2.73 & 2.34 & 1.34 & 2.11 & 2.49 & 2.3 & 1.97 & 2.73 & 1.94 & 2.36 & 1.95 & 2.32 & 2.65 & 3.64 \\ 
   & SF & 20.72 & 27.71 & 8.54 & 20.18 & 23.06 & 15.35 & 14.27 & 13.05 & 9.37 & 27.94 & 40.03 & 16.3 & 20.81 & 6.78 \\ 
   & baseline & 22.76 & 5.34 & 5.98 & 5.84 & 5.16 & 10.37 & 10.9 & 7.26 & 10.25 & 10.18 & 6.26 & 5.69 & 6.71 & 10.31 \\ 
   & warm & 23.46 & 7.2 & 5.75 & 6.41 & 5.21 & 11.56 & 12.21 & 8.77 & 8.93 & 6.81 & 8.93 & 6.15 & 7.25 & 18.98 \\ 
   \hline
{\textbf{X-NN}} & frozen & 6.63 & 14.04 & 10.73 & 19.57 & 17.94 & 14 & 13.67 & 18.83 & 14.36 & 10.81 & 12.25 & 16.61 & 39.96 & 32.81 \\ 
   & multi head & 1.19 & 11.38 & 84.13 & 174.18 & 1.87 & 19.55 & 62.12 & 22.7 & 9.67 & {\textbf{0.85*}} & 3.34 & 5.03 & {\textbf{0.94*}} & 111.42 \\ 
   & SF & 18.83 & 10.72 & 10.3 & 10.61 & 10.11 & 12.5 & 21.37 & 11.12 & 8.27 & 16.33 & 9.81 & 13.8 & 9.85 & 10.15 \\ 
   & baseline & 19.52 & 8.25 & 4.68 & 5.06 & 6.6 & 11.5 & 10.78 & 12 & 10.26 & 6.25 & 13.11 & 7.27 & 9.2 & 18.45 \\ 
   & warm & 20.06 & 8.54 & 5.57 & 6.37 & 10.77 & 9.34 & 11.36 & 13.16 & 8.03 & 8.66 & 10.96 & 7.52 & 11.57 & 16.7 \\ 
   \hline
{\textbf{Y-NN}} & frozen & 1.54 & {\textbf{1*}} & 2.1 & 3.3 & {\textbf{1.11*}} & 46.07 & 27.63 & 5.47 & 7.48 & 7.21 & 1.02 & {\textbf{1.15*}} & 0.97 & 43.82 \\ 
   & multi head & 0.92 & 2.21 & 1.26 & 15.86 & 1.19 & 20.47 & 1.76 & 2.75 & 3.96 & 9.4 & 19.11 & 1.26 & 35.43 & 9.29 \\ 
   & SF & 5.68 & 12 & 16.57 & 38 & 5.54 & 3.49 & 8.2 & 4.8 & 2.61 & 7.21 & 8.48 & 12.98 & 9.58 & 5.08 \\ 
   & baseline & {\textbf{0.9*}} & 1.31 & 5.24 & 31.43 & 6.8 & {\textbf{1.23*}} & 7.5 & {\textbf{1.07*}} & 1.32 & 1.35 & 8.19 & 1.2 & 4.11 & 7.72 \\ 
   & warm & 48.91 & 1.26 & 5.79 & 3.61 & 29.43 & 2.71 & 3.94 & 12.87 & 21.76 & 13.25 & 15.78 & 17.45 & 1.12 & 29.01 \\ 
   \hline
{\textbf{joint}} & joint & 13.75 & 5.81 & 3.46 & 4.12 & 3.46 & 12.22 & 9.58 & 14.96 &  & 4.75 & 8.55 & 13.13 & 9.91 & 11.68 \\ 
   \hline
{\textbf{MLRW}} &  & 1 & 1.41 & {\textbf{1.05*}} & {\textbf{1.94*}} & 1.97 & 1.26 & {\textbf{1.03*}} & 2.05 & {\textbf{0.9*}} & 1.85 & {\textbf{1*}} & 1.15 & 1.57 & {\textbf{2.75*}} \\ 
  \end{tabular}
\caption{MSE in percent for different CATE estimators.} 
\label{table:LM}
\end{flushleft}
\end{sidewaystable}

% latex table generated in R 3.4.4 by xtable 1.8-2 package
% Thu May 17 18:43:51 2018
\begin{sidewaystable}[ht]
\begin{flushleft}
\begin{tabular}{rrrrrrrrrrrrrrrr}
 {} & {\textbf{Method}} & {\textbf{RF-1}} & {\textbf{RF-2}} & {\textbf{RF-3}} & {\textbf{RF-4}} & {\textbf{RF-5}} & {\textbf{RF-6}} & {\textbf{RF-7}} & {\textbf{RF-8}} & {\textbf{RF-9}} & {\textbf{RF-10}} & {\textbf{RF-11}} & {\textbf{RF-12}} & {\textbf{RF-13}} & {\textbf{RF-14}} \\ 
  \hline
{\textbf{t-lm}} &  & 17.95 & 26.73 & 2.48 &  & 4.84 & 1.76 & 6.2 & 7.55 & 24.58 & 3.5 & 2.35 & 2.1 & 3.39 & 7.41 \\ 
   \hline
{\textbf{s-rf}} &  & 7.18 & 7 & 4.18 & 6.86 & 8.43 & 9.15 & 6.03 & 9.28 & 5.95 & 8.49 & 6.7 & 8.1 & 6.46 & 8.64 \\ 
   \hline
{\textbf{t-rf}} &  & 10.95 & 7.76 & 4.69 & 7.79 & 9.03 & 10.07 & 6.67 & 10.08 & 6.3 & 10.35 & 8.19 & 9.28 & 6.73 & 10.04 \\ 
   \hline
{\textbf{R-NN}} &  & 11.41 & 1.18 & 5.7 & 1.26 & 0.7 & 3.79 & 9.11 & 8.01 & 7.17 & 5.14 & 4.15 & 0.52 & 1.26 & 6.88 \\ 
   \hline
{\textbf{S-NN}} & frozen & 1.56 & 0.66 & 0.7 & 0.77 & 0.87 & 0.83 & 0.91 & 0.89 & 0.59 & 0.87 & 0.6 & 0.92 & 0.91 & 1.23 \\ 
   & multi head & 0.69 & 1.23 & 0.4 & 1.59 & 1.58 & 0.65 & 0.81 & 0.57 & {\textbf{0.28*}} & 1.98 & 1.08 & 1.24 & 2.27 & 3.66 \\ 
   & SF & 1.36 & 9.64 & 4.26 & 4.53 & 6 & 6.74 & 11.68 & 2.8 & 12.01 & 41.6 & 45.72 & 4.59 & 1.98 & 6.75 \\ 
   & baseline & 0.91 & 1.28 & 0.84 & 0.93 & 1.78 & 2.05 & 1.16 & 2.04 & 1.61 & 1.06 & 1.29 & 1.49 & 1.99 & 2.69 \\ 
   & warm & 1.08 & 1.3 & 0.94 & 1.16 & 1.85 & 2.22 & 1.4 & 2.15 & 1.65 & 1.12 & 1.37 & 1.42 & 1.83 & 2.52 \\ 
   \hline
{\textbf{T-NN}} & frozen & 1.55 & 1.02 & 0.62 & 0.95 & 1.11 & 0.99 & 0.89 & 1.18 & 0.86 & 1.03 & 0.89 & 1.01 & 1.14 & 1.49 \\ 
   & multi head & 0.66 & 0.58 & 0.35 & 0.53 & 0.63 & 0.53 & 0.47 & 0.63 & 0.49 & 0.57 & 0.51 & 0.55 & 0.64 & 0.91 \\ 
   & SF & 11.66 & 27.5 & 6.89 & 22.59 & 20.99 & 12.04 & 19.77 & 7.85 & 5.39 & 17.31 & 17.15 & 9.1 & 30.42 & 3.55 \\ 
   & baseline & 7.83 & 4.25 & 1.44 & 1.47 & 1.35 & 1.33 & 2.05 & 8.55 & 1.79 & 3.21 & 2.29 & 3.62 & 2.03 & 8.7 \\ 
   & warm & 12.74 & 3.32 & 1.33 & 1.58 & 1.06 & 1.6 & 2.19 & 9.28 & 1.52 & 2.77 & 3.73 & 4.99 & 1.64 & 8.74 \\ 
   \hline
{\textbf{X-NN}} & frozen & 2.53 & 22.89 & 55.57 & 44.11 & 4.18 & 5.72 & 33.18 & 2.62 & 7.59 & 4.96 & 4.41 & 2.01 & 3.01 & 1.23 \\ 
   & multi head & 3.45 & 2.53 & 47.09 & 39.7 & 2 & 27.62 & 10.39 & 0.78 & 11.8 & 10.85 & 0.93 & 9.72 & 11.74 & 30.62 \\ 
   & SF & 4.34 & 1.85 & 5.47 & 6.82 & 4.21 & 11.89 & 7.91 & 5.02 & 4.67 & 6.82 & 16.39 & 6.99 & 5.22 & 2.11 \\ 
   & baseline & 4.09 & 2.05 & 1.67 & 0.73 & 0.58 & 1.16 & 4.97 & 9.23 & 4.05 & 1.49 & 5.17 & 2.15 & 4.07 & 8.05 \\ 
   & warm & 2.51 & 1.73 & 1.72 & 1.48 & 0.97 & 1.04 & 7.81 & 4.35 & 4.59 & 1.34 & 3.86 & 1.37 & 2.08 & 4.83 \\ 
   \hline
{\textbf{Y-NN}} & frozen & 0.42 & {\textbf{0.32*}} & 10.65 & 0.72 & 36 & 1.62 & 0.82 & 0.87 & 15.19 & 1.61 & 0.77 & 46.3 & 28.37 & 1.98 \\ 
   & multi head & 29.41 & 9.71 & 2.74 & 12.27 & 48.87 & 16.59 & 50.21 & 73.74 & 1.65 & 1.76 & 3.52 & 24.26 & 0.76 & 216.69 \\ 
   & SF & 5.45 & 2.51 & 1.53 & 4.76 & 9.61 & 20.81 & 3.68 & 4.84 & 10.8 & 2.87 & 2.78 & 1.36 & 9.07 & 4.08 \\ 
   & baseline & 0.71 & 1.06 & 1.54 & 0.57 & 2.44 & 0.47 & 0.73 & 1.38 & 1.25 & 1.77 & 0.71 & 0.29 & 0.63 & 1.98 \\ 
   & warm & 0.9 & 27.01 & 0.52 & 22.7 & 1.37 & 24.94 & {\textbf{0.3*}} & 2.08 & 12.36 & 3.39 & 9.58 & 2.74 & 2.47 & 11.48 \\ 
   \hline
{\textbf{joint}} & joint &  & 33.47 & 1.3 & 0.61 & 5.15 & 0.47 & 2.68 &  & 2.67 & {\textbf{0.37*}} & 7.76 & 28.14 & 7.48 & 10.04 \\ 
   \hline
{\textbf{MLRW}} &  & {\textbf{0.37*}} & 1.12 & {\textbf{0.18*}} & {\textbf{0.3*}} & {\textbf{0.46*}} & {\textbf{0.21*}} & 0.42 & {\textbf{0.45*}} & 0.37 & 1.18 & {\textbf{0.35*}} & {\textbf{0.23*}} & 0.26 & {\textbf{0.16*}} \\ 
  \end{tabular}
\caption{MSE in percent for different CATE estimators.} 
\label{table:RF}
\end{flushleft}
\end{sidewaystable}

% latex table generated in R 3.4.4 by xtable 1.8-2 package
% Thu May 17 18:43:51 2018
\begin{sidewaystable}[ht]
\begin{flushleft}
\begin{tabular}{rrrrrrrrrrrrrrrr}
 {} & {\textbf{Method}} & {\textbf{RFt-1}} & {\textbf{RFt-2}} & {\textbf{RFt-3}} & {\textbf{RFt-4}} & {\textbf{RFt-5}} & {\textbf{RFt-6}} & {\textbf{RFt-7}} & {\textbf{RFt-8}} & {\textbf{RFt-9}} & {\textbf{RFt-10}} & {\textbf{RFt-11}} & {\textbf{RFt-12}} & {\textbf{RFt-13}} & {\textbf{RFt-14}} \\ 
  \hline
{\textbf{t-lm}} &  & 22.65 & 2.74 & 2.3 & 5.07 & 4.78 & 8.3 & 9.77 & 38.73 & 23.94 & 8.94 & 83.37 & 4.26 & 4.31 & 31.76 \\ 
   \hline
{\textbf{s-rf}} &  & 3.54 & 7.04 & 4.39 & 6.4 & 8.09 & 6.41 & 11.44 & 6.33 & 6.05 & 6.52 & 5.09 & 9.29 & 11.64 & 4.28 \\ 
   \hline
{\textbf{t-rf}} &  & 13.3 & 13.55 & 8.41 & 11.9 & 14.11 & 15.48 & 14.42 & 12.54 & 10.19 & 18.99 & 12.48 & 14.64 & 13.61 & 12.02 \\ 
   \hline
{\textbf{R-NN}} &  & 62.08 & 17.68 & 5.22 & 19.91 & 14.43 & 20.53 & 18.51 & 49.63 & 8.99 & 86.89 & 54.17 & 23.86 & 3.02 & 18.98 \\ 
   \hline
{\textbf{S-NN}} & frozen & 2.49 & 51.39 & 37.47 & 49.5 & 43.17 & 29.19 & 21.81 & 27.25 & 61.34 & 27.02 & 17.25 & 64.22 & 24.94 & 11.24 \\ 
   & multi head & 1.4 & 0.57 & 0.73 & 1.48 & 1.54 & 1.27 & 1.1 & 1.63 & 0.88 & 1.07 & 1.53 & 0.79 & 0.74 & 1.59 \\ 
   & SF & 0.84 & 68.31 & 4.12 & 10.77 & 17.28 & 20.06 & 10.22 & 1.08 & 8.95 & 9.11 & 5.97 & 19.67 & 11.64 & 8.68 \\ 
   & baseline & 14.66 & 2.08 & 0.95 & 1.38 & 2.24 & 7.3 & 3.17 & 1.08 & 0.88 & 5.29 & 1.44 & 2.37 & 2.85 & 1.78 \\ 
   & warm & 17.64 & 2.76 & 1.38 & 2.31 & 1.99 & 8.64 & 4.32 & 1.37 & 0.85 & 5.5 & 3.76 & 5.44 & 3.13 & 3.49 \\ 
   \hline
{\textbf{T-NN}} & frozen & 2.57 & 1.69 & 1.12 & 1.55 & 1.81 & 1.6 & 1.52 & 1.86 & 1.43 & 1.63 & 1.47 & 1.68 & 1.85 & 2.27 \\ 
   & multi head & 0.69 & 0.55 & 0.34 & 0.49 & 0.63 & 0.64 & 0.51 & 0.64 & 0.5 & 0.53 & {\textbf{0.46*}} & {\textbf{0.55*}} & 0.66 & 0.8 \\ 
   & SF & 18.04 & 9.54 & 7.33 & 6.61 & 22.11 & 9.21 & 20.2 & 10.65 & 15.06 & 13.87 & 27.64 & 17.96 & 14.08 & 4.49 \\ 
   & baseline & 43.37 & 14.55 & 9.48 & 19.35 & 15.68 & 20.91 & 16 & 41.72 & 6.99 & 30.02 & 45.59 & 16.03 & 4.79 & 18.63 \\ 
   & warm & 69.41 & 19.81 & 12.99 & 24.12 & 20.28 & 37.07 & 19.64 & 39.37 & 9.36 & 40.65 & 37.59 & 18.62 & 6.87 & 32.52 \\ 
   \hline
{\textbf{X-NN}} & frozen & 12.41 & 221.49 & 215.01 & 131.46 & 94.89 & 328.53 & 199.48 & 41.14 & 398.33 & 83.75 & 132 & 434.25 & 125.88 & 168.92 \\ 
   & multi head & 12.24 & 0.79 & 3.22 & 12.46 & 2.87 & 12.11 & 21.38 & 1.86 & 1.48 & 18.46 & 1.82 & 99.38 & 14.67 & 28.43 \\ 
   & SF & 3.98 & 11.03 & 4.33 & 5.31 & 5.42 & 4.67 & 14.03 & 5.99 & 7 & 7 & 6.42 & 7.21 & 4.89 & 2.34 \\ 
   & baseline & 60.44 & 15.03 & 4.93 & 7.08 & 9.51 & 21.22 & 14.7 & 16.83 & 8 & 80.71 & 23.89 & 25.5 & 4.67 & 17.1 \\ 
   & warm & 30.46 & 13.02 & 5.11 & 9.95 & 9.21 & 13.12 & 9.95 & 12.8 & 4.48 & 50.45 & 19.27 & 12.92 & 8.1 & 15.73 \\ 
   \hline
{\textbf{Y-NN}} & frozen & 1.72 & {\textbf{0.33*}} & 0.87 & 2.47 & 1.8 & 5.44 & 1.35 & 19.37 & 37.98 & 0.81 & 0.7 & 33.11 & 1.34 & 13.64 \\ 
   & multi head & 5.95 & 11.12 & 0.41 & 2.94 & 2.49 & 24.73 & 3.44 & 58.63 & 12.99 & 0.41 & 9.55 & 0.6 & 3.93 & 0.49 \\ 
   & SF & 3.98 & 14.56 & 12.85 & 4.3 & 3.66 & 9.78 & 11.13 & 3.03 & 5.78 & 3.37 & 11.87 & 11.33 & 4.51 & 2.64 \\ 
   & baseline & 0.88 & 10.92 & 0.28 & 1.32 & 1.89 & 11.35 & 8.42 & 1.6 & 5.66 & 0.54 & 0.54 & 2.01 & 0.61 & 2.13 \\ 
   & warm & 8.03 & 0.48 & 0.36 & 0.71 & 2.53 & 1.21 & {\textbf{0.31*}} & 1.77 & {\textbf{0.41*}} & 0.74 & 23.48 & 1.4 & 1.15 & 17.72 \\ 
   \hline
{\textbf{joint}} & joint & 102.83 & 38.47 & 133.88 & 15.31 & 37.32 & 410.04 & 38.28 & 14.96 & 5.02 & 11.97 & 33.71 & 26.33 & 11.99 & 147.45 \\ 
   \hline
{\textbf{MLRW}} &  & {\textbf{0.31*}} & 1.25 & {\textbf{0.21*}} & {\textbf{0.27*}} & {\textbf{0.35*}} & {\textbf{0.29*}} & 0.33 & {\textbf{0.61*}} & 0.61 & {\textbf{0.4*}} & 0.57 & 1.32 & 0.8 & {\textbf{0.24*}} \\ 
  \end{tabular}
\caption{MSE in percent for different CATE estimators.} 
\label{table:RF1}
\end{flushleft}
\end{sidewaystable}




	
\end{document}
